% 此文件由本卷已核验事件账本生成；仅作附录跳转导航。
\section{完整事件导航}
本节为账本中尚未在正文首次出现的事件提供跳转目标。它保留日期、政权和事件名称，供索引、读者与后续增写定位；它不是正文，也不以一行标题替代事件的原因、选择、后果和史料讨论。每条的核心来源、可信度与争议说明见事件账本及\hyperref[app:sources]{来源附录}。
\begin{description}
  \item[太熙元年四月（290年）] \eventanchor{JH-0290-WU-EMPEROR-DEATH}{西晋·司马炎去世，惠帝司马衷即位，杨骏以太傅录尚书事主持朝政}
  \item[永平元年（291年）] \eventanchor{JH-0291-JIA-OVERTHROWS-YANG}{西晋·贾后联络楚王司马玮等发动政变，杨骏被杀，杨氏外戚集团被清除}
  \item[元康元年（291年）] \eventanchor{JH-0291-LIANG-WEI-PURGE}{西晋·汝南王司马亮与卫瓘辅政后又被贾后借楚王司马玮诛杀，司马玮旋即被处死}
  \item[元康六年至九年（296年至299年）] \eventanchor{JH-0296-QI-WANNIAN}{西晋与关陇氐羌集团·氐族齐万年在关中起兵，西晋多次征讨并长期调动关陇军粮，至299年叛乱被平定}
  \item[永康元年（300年）] \eventanchor{JH-0300-HEIR-APPARENT-KILLED}{西晋·贾后以谋反罪废杀太子司马遹，皇位继承与宫廷联盟由此公开破裂}
  \item[永宁元年（301年）] \eventanchor{JH-0301-SIMA-LUN-USURPATION}{西晋·赵王司马伦杀贾后并废惠帝自立为帝，短暂改变西晋帝统}
  \item[永宁元年（301年）] \eventanchor{JH-0301-SIMA-JIONG-COUP}{西晋·齐王司马冏等起兵推翻司马伦，复立惠帝并以大司马辅政}
  \item[太安元年（302年）] \eventanchor{JH-0302-SIMA-YING-DEFEATS-JIONG}{西晋·河间王司马颙与成都王司马颖起兵，长沙王司马乂在洛阳杀齐王司马冏并接掌朝政}
  \item[太安二年（303年）] \eventanchor{JH-0303-SIMA-AI-DEFENDS-LUOYANG}{西晋·长沙王司马乂据洛阳抵抗河间王司马颙和成都王司马颖联军，京师陷入长期战斗}
  \item[太安二年（303年）] \eventanchor{JH-0303-LI-XIONG-TAKES-CHENGDU}{成汉与益州西晋守军·巴氐首领李雄攻取成都，益州的西晋统治被成汉政权取代}
  \item[永安元年（304年）] \eventanchor{JH-0304-SIMA-AI-KILLED}{西晋·司马乂被东海王司马越等交给河间王司马颙军并遭杀害，洛阳政权落入外军支配}
  \item[永兴元年（304年）] \eventanchor{JH-0304-LIU-YUAN-HAN}{汉赵与并州诸部·刘渊在左国城起兵，自称汉王，汇聚匈奴五部及并州反晋力量}
  \item[永兴元年（304年）] \eventanchor{JH-0304-CHEN-MIN-REBELLION}{西晋与江东地方集团·陈敏据江东起兵并自置官属，试图利用八王之乱切割东南地区}
  \item[永兴二年（305年）] \eventanchor{JH-0305-SIMA-YUE-COALITION}{西晋·东海王司马越以讨河间王司马颙为名起兵，八王之乱扩大为跨州郡的长期战争}
  \item[光熙元年（306年）] \eventanchor{JH-0306-HUI-EMPEROR-DEATH}{西晋·惠帝在洛阳去世，司马炽即位为怀帝，司马越继续掌握主要军政权}
  \item[晏平元年（306年）] \eventanchor{JH-0306-LI-XIONG-EMPEROR}{成汉·李雄在成都称帝，国号成，使巴蜀割据由起事集团转为帝国体制}
  \item[永凤元年（308年）] \eventanchor{JH-0308-LIU-YUAN-EMPEROR}{汉赵·刘渊称帝，正式以汉为国号并以平阳为政治中心}
  \item[永嘉三年（309年）] \eventanchor{JH-0309-HAN-ZHAO-SIEGE}{汉赵与西晋·刘聪、王弥等围攻洛阳，西晋京师依靠临时征兵和城防勉强守住}
  \item[河瑞二年（310年）] \eventanchor{JH-0310-LIU-CONG-ACCESSION}{汉赵·刘渊去世后刘聪经宫廷斗争继位，汉赵继续对西晋发动攻势}
  \item[永嘉四年（310年）] \eventanchor{JH-0310-SIMA-YUE-DEATH}{西晋·司马越在东出避乱途中去世，其部众和辎重随即失去统一指挥}
  \item[永嘉五年六月（311年）] \eventanchor{JH-0311-YONGJIA-FALL}{汉赵与西晋·汉赵军攻陷洛阳，大量官民死散，西晋都城体系遭到毁灭性破坏}
  \item[永嘉五年（311年）] \eventanchor{JH-0311-HUAI-EMPEROR-CAPTURED}{汉赵与西晋·怀帝在洛阳陷落时被汉赵俘获并押往平阳，西晋皇统陷入屈辱性危机}
  \item[建兴元年（313年）] \eventanchor{JH-0313-HUAI-EMPEROR-KILLED}{汉赵与西晋·汉赵杀害被俘的怀帝，长安的司马邺被拥立为愍帝}
  \item[建兴元年（313年）] \eventanchor{JH-0313-ZU-TI-NORTHERN-EXPEDITION}{东晋前身与河南地方集团·祖逖率部北渡长江，经屯田与结盟经营豫州、河南一带}
  \item[建兴三年（315年）] \eventanchor{JH-0315-TUOBA-YILU}{代国与西晋残余政权·拓跋猗卢受晋封为代王，在北部边地扩展部众与城郭控制}
  \item[建兴四年十一月（316年）] \eventanchor{JH-0316-CHANGAN-FALL}{汉赵与西晋·刘曜攻陷长安，愍帝出降，西晋在北方的最后朝廷覆灭}
  \item[建武元年三月（317年）] \eventanchor{JH-0317-SIMA-RUI-KING-JIN}{东晋·司马睿在建康即晋王位，建立承继西晋的南方朝廷框架}
  \item[太兴二年（319年）] \eventanchor{JH-0319-SHI-LE-LATER-ZHAO}{后赵·石勒在襄国称赵王并建立独立政权，后赵由汉赵将领集团分出}
  \item[光初二年（319年）] \eventanchor{JH-0319-LIU-YAO-RENAMES-ZHAO}{前赵·刘曜在长安改汉国号为赵，以关中为中心重建政权身份}
  \item[大兴三年（320年）] \eventanchor{JH-0320-DU-TAO-DEFEATED}{东晋与荆湘地方集团·杜弢领导的荆湘叛乱被王敦等镇压，东晋重新取得长江中游部分军政控制}
  \item[大兴四年（321年）] \eventanchor{JH-0321-WANG-DUN-FIRST-REBELLION}{东晋·王敦以清君侧为名自武昌东下，迫使元帝调整中枢人事并承认其军事优势}
  \item[永昌元年（322年）] \eventanchor{JH-0322-WANG-DUN-TAKES-JIANKANG}{东晋·王敦第二次起兵攻入建康，杀逐元帝近臣并控制朝政}
  \item[永昌元年闰十一月（323年）] \eventanchor{JH-0323-YUAN-DEATH}{东晋·元帝去世，明帝司马绍即位，王敦集团与建康朝廷的对抗延续}
  \item[太宁二年（324年）] \eventanchor{JH-0324-WANG-DUN-COLLAPSE}{东晋·王敦再度举兵但病死于途中，其部众在建康内外被明帝集团击败}
  \item[太宁三年（325年）] \eventanchor{JH-0325-ZU-TI-DEATH}{东晋与河南地方集团·祖逖去世，其北伐军和河南经营失去核心协调者}
  \item[咸和二年（327年）] \eventanchor{JH-0327-SU-JUN-REBELLION}{东晋·历阳将苏峻以拒绝征召为由起兵，联合祖约威胁建康}
  \item[咸和三年（328年）] \eventanchor{JH-0328-SU-JUN-TAKES-JIANKANG}{东晋·苏峻军攻入建康，劫持成帝并控制宫城和朝廷}
  \item[太和三年（328年）] \eventanchor{JH-0328-SHI-LE-CAPTURES-LIU-YAO}{后赵与前赵·石勒在洛阳附近击败并俘获前赵皇帝刘曜，前赵主力由此崩解}
  \item[太和四年（329年）] \eventanchor{JH-0329-LATER-ZHAO-DESTROYS-FORMER-ZHAO}{后赵与前赵·石勒攻取长安并处置刘曜之子刘熙，前赵政权灭亡}
  \item[建平元年（330年）] \eventanchor{JH-0330-SHI-LE-EMPEROR}{后赵·石勒在襄国称帝，后赵以河北为核心建立更完整的官僚与军政体系}
  \item[建平四年（333年）] \eventanchor{JH-0333-SHI-LE-DEATH}{后赵·石勒去世，太子石弘继位，石虎掌握军政实权}
  \item[建武元年（334年）] \eventanchor{JH-0334-SHI-HU-COUP}{后赵·石虎废杀石弘，自立为天王，后赵皇位转入其支系}
  \item[玉衡二十四年（334年）] \eventanchor{JH-0334-LI-XIONG-DEATH}{成汉·李雄去世，李班继位后旋即被李期夺权，成汉内部继承冲突公开化}
  \item[咸康元年（335年）] \eventanchor{JH-0335-MURONG-HUANG-YAN}{前燕·慕容皝在辽西称燕王，前燕政权以棘城为中心扩展}
  \item[咸康二年（336年）] \eventanchor{JH-0336-MURONG-HUANG-DEFEATS-DUAN}{前燕与段部·慕容皝击败段辽并整合辽西部众，前燕的军事与人口资源扩大}
  \item[咸康四年（338年）] \eventanchor{JH-0338-SHI-HU-ATTACKS-YAN}{后赵与前燕·石虎大军进攻前燕，慕容皝依托棘城及周边地形击退来军}
  \item[咸康四年（338年）] \eventanchor{JH-0338-LI-SHOU-USURPS-CHENGHAN}{成汉·李寿废李期自立并改国号为汉，成汉政权的宗室结构再次重组}
  \item[咸康七年（341年）] \eventanchor{JH-0341-WANG-XINGZHI-EPITAPH}{东晋江南·王兴之夫妇墓志下葬，为建康周边侨姓官员家庭、婚姻与官职提供有纪年实物材料}
  \item[咸康八年（342年）] \eventanchor{JH-0342-MURONG-HUANG-GOGURYEO}{前燕与高句丽·慕容皝进攻高句丽都城丸都，掳掠人口并改变辽东力量关系}
  \item[永和二年（346年）] \eventanchor{JH-0346-HUAN-WEN-INVADES-CHENGHAN}{东晋与成汉·桓温受命自荆州西征成汉，沿长江上游进入巴蜀}
  \item[永和三年（347年）] \eventanchor{JH-0347-CHENGHAN-FALL}{东晋与成汉·桓温攻入成都，李势投降，成汉灭亡，东晋恢复益州控制}
  \item[永宁元年（349年）] \eventanchor{JH-0349-SHI-HU-DEATH}{后赵·石虎去世，诸子与养孙冉闵围绕都城邺城和军队展开继承争夺}
  \item[永宁二年（350年）] \eventanchor{JH-0350-RAN-MIN-WEI}{冉魏与后赵遗众·冉闵废石祗后在邺城称帝，建立冉魏并在内战中针对羯胡及后赵旧部实施暴力清洗}
  \item[永宁二年（350年）] \eventanchor{JH-0350-FU-HONG-CLAIMS-QIN}{前秦前身·氐族首领苻洪在关中东部聚众，自称三秦王，前秦政权的军事基础开始形成}
  \item[永宁三年（351年）] \eventanchor{JH-0351-FU-JIAN-FOUNDS-FORMER-QIN}{前秦·苻健在长安称天王，建立前秦并开始重建关中官署}
  \item[永和八年（352年）] \eventanchor{JH-0352-RAN-MIN-DEFEATED}{冉魏与前燕·冉闵在魏昌一带败于前燕，被俘后处死，冉魏政权迅速瓦解}
  \item[永和九年三月三日（353年）] \eventanchor{JH-0353-LANTING-GATHERING}{东晋·王羲之与会稽士人在兰亭修禊并作序，成为东晋士人交游与书写文化的标志性事件}
  \item[永和十年（354年）] \eventanchor{JH-0354-HUAN-WEN-NORTHERN-EXPEDITION}{东晋与前秦·桓温自荆州北伐前秦，进至关中外围后因粮运困难撤退}
  \item[永和十一年（355年）] \eventanchor{JH-0355-LATER-ZHAO-ENDS}{后赵残余与冉魏·石祗在襄国被杀，后赵最后的政治中心消失}
  \item[永和十二年（356年）] \eventanchor{JH-0356-HUAN-WEN-RECOVERS-LUOYANG}{东晋与前燕·桓温北伐进入洛阳，修复晋帝陵并短暂恢复东晋在中原的象征性存在}
  \item[永和十三年（357年）] \eventanchor{JH-0357-FU-JIAN-ACCESSION}{前秦·苻坚废杀苻生即位，以王猛等整顿前秦朝政}
  \item[永和十三年（357年）] \eventanchor{JH-0357-WANG-MENG-REFORMS}{前秦·苻坚任用王猛参与中枢政务，以整饬吏治和抑制豪强为前秦集权奠基}
  \item[升平四年（360年）] \eventanchor{JH-0360-MURONG-JUN-EMPEROR}{前燕·慕容儁在邺称帝，前燕由东北政权转为据有中原部分地区的帝国}
  \item[升平五年（361年）] \eventanchor{JH-0361-AI-EMPEROR-ACCESSION}{东晋·穆帝去世，哀帝司马丕即位，东晋朝政继续由会稽王司马昱等辅政集团主导}
  \item[隆和元年（362年）] \eventanchor{JH-0362-MURONG-JUN-DEATH}{前燕·慕容儁去世，年幼的慕容暐继位，慕容恪等辅政}
  \item[兴宁三年（365年）] \eventanchor{JH-0365-FORMER-YAN-TAKES-LUOYANG}{前燕与东晋·慕容恪攻取洛阳，前燕控制中原旧都并迫使东晋北方据点后撤}
  \item[太和二年（367年）] \eventanchor{JH-0367-FORMER-QIN-REBELLIONS}{前秦·苻双、苻柳等宗室和地方势力反叛，苻坚与王猛平定内乱}
  \item[太和四年（369年）] \eventanchor{JH-0369-FANGTOU-DEFEAT}{东晋与前燕·桓温北伐前燕，在枋头受慕容垂等反击而败，退军途中损失严重}
  \item[太和五年（370年）] \eventanchor{JH-0370-FORMER-QIN-CONQUERS-YAN}{前秦与前燕·王猛率前秦军攻灭前燕，慕容暐降服，前秦取得河北和中原北部}
  \item[咸安元年（371年）] \eventanchor{JH-0371-FORMER-QIN-ANNEXES-CHOUCHI}{前秦与仇池氐族政权·前秦攻取仇池，控制陇南山地通道并扩展对氐族地区的影响}
  \item[咸安二年（372年）] \eventanchor{JH-0372-XIAOWU-EMPEROR}{东晋·简文帝去世，孝武帝司马曜即位，谢安等高门重臣参与辅政}
  \item[宁康元年（373年）] \eventanchor{JH-0373-FORMER-QIN-TAKES-LIANG-YI}{前秦与东晋·前秦军攻取东晋梁、益二州北部州郡，深入汉中与四川北缘}
  \item[太元元年（376年）] \eventanchor{JH-0376-FORMER-QIN-DESTROYS-LIANG}{前秦与前凉·前秦攻灭前凉，张氏河西政权终结，凉州纳入前秦势力}
  \item[太元元年（376年）] \eventanchor{JH-0376-FORMER-QIN-DESTROYS-DAI}{前秦与代国·前秦分兵攻灭代国，拓跋什翼犍败亡，拓跋部暂时分散}
  \item[太元三年（378年）] \eventanchor{JH-0378-XIANGYANG-SIEGE}{前秦与东晋·苻丕等围攻襄阳，前秦以汉水交通枢纽为目标向东晋施压}
  \item[太元三年（378年）] \eventanchor{JH-0378-BEIFU-ARMY}{东晋·谢玄在京口训练北府兵，东晋以侨民和北来军人为核心重组淮南防务}
  \item[太元四年（379年）] \eventanchor{JH-0379-XIANGYANG-FALL}{前秦与东晋·襄阳城陷，守将朱序被俘，前秦控制汉水中游重镇}
  \item[太元七年（382年）] \eventanchor{JH-0382-FU-JIAN-SOUTHERN-INVASION}{前秦与东晋·苻坚决定大举南伐东晋，调集北方多地兵员与粮运资源}
  \item[太元九年（384年）] \eventanchor{JH-0384-MURONG-CHUI-LATER-YAN}{后燕·慕容垂在河北起兵，建立后燕并号召慕容旧部反前秦}
  \item[太元九年（384年）] \eventanchor{JH-0384-YAO-CHANG-LATER-QIN}{后秦·姚苌在北地起兵并称万年秦王，后秦政权自前秦军中分出}
  \item[太元十年（385年）] \eventanchor{JH-0385-FU-JIAN-KILLED}{前秦与后秦·苻坚在五将山被姚苌部将杀害，前秦失去核心统治者}
  \item[太元十一年（386年）] \eventanchor{JH-0386-TUOBA-GUI-RESTORES-WEI}{北魏前身·拓跋珪在牛川重建拓跋政权并称代王，旋改国号魏}
  \item[太元十一年（386年）] \eventanchor{JH-0386-LU-GUANG-LATER-LIANG}{后凉·吕光自西域东归后据凉州建立后凉，河西走廊脱离前秦控制}
  \item[太元十一年（386年）] \eventanchor{JH-0386-WESTERN-YAN}{西燕·慕容泓、慕容永等慕容部众在关中、河东活动并建立西燕政权}
  \item[太元十七年（392年）] \eventanchor{JH-0392-MURONG-CHUI-DEATH}{后燕·慕容垂去世，慕容宝继位，后燕在北魏与内部矛盾压力下转入不稳}
  \item[太元十八年（393年）] \eventanchor{JH-0393-YAO-XING-ACCESSION}{后秦·姚苌去世，姚兴继位，后秦继续整合关中并调整对前秦残余的战争}
  \item[太元十九年（394年）] \eventanchor{JH-0394-LATER-QIN-DESTROYS-FORMER-QIN}{后秦与前秦残余·姚兴军击败并杀害苻登，前秦残余政权灭亡}
  \item[太元二十年（395年）] \eventanchor{JH-0395-CANHEBEI}{北魏与后燕·拓跋珪在参合陂大败后燕军，北魏取得对河北竞争的主动}
  \item[皇始元年（396年）] \eventanchor{JH-0396-TUOBA-GUI-EMPEROR}{北魏·拓跋珪称帝，北魏以平城为核心扩张并向后燕发动攻势}
  \item[皇始二年（397年）] \eventanchor{JH-0397-NORTHERN-WEI-TAKES-ZHONGSHAN}{北魏与后燕·北魏攻取中山，后燕政权被迫向东北退却，河北主导权易手}
  \item[天兴元年（398年）] \eventanchor{JH-0398-NORTHERN-WEI-MOVES-PINGCHENG}{北魏·拓跋珪迁都平城并重建都城机构，北魏的统治重心由代北部落空间转向城郭行政}
  \item[隆安三年（399年）] \eventanchor{JH-0399-SUN-EN-REBELLION}{东晋与浙东天师道集团·孙恩在会稽起兵，东晋东南沿海州郡陷入反复征战}
  \item[弘始三年（401年）] \eventanchor{JH-0401-KUMARAJIVA-CHANGAN}{后秦·鸠摩罗什抵达长安，姚兴支持译经活动，后秦成为重要佛教翻译中心}
  \item[元兴元年（402年）] \eventanchor{JH-0402-HUAN-XUAN-TAKES-JIANKANG}{东晋·桓玄自荆州东下攻入建康，控制安帝和东晋朝政}
  \item[元兴二年十二月（403年）] \eventanchor{JH-0403-HUAN-XUAN-USURPATION}{桓楚与东晋·桓玄废安帝自立为楚帝，东晋皇统短暂中断}
  \item[元兴三年（404年）] \eventanchor{JH-0404-LIU-YU-RESTORES-JIN}{东晋·刘裕、刘毅等起兵推翻桓楚，迎复安帝，东晋皇统恢复}
  \item[义熙元年（405年）] \eventanchor{JH-0405-CUAN-BAOZI-STELE}{东晋宁州地方社会·《爨宝子碑》立于宁州，记录爨氏地方官员身份与家族记忆}
  \item[义熙二年（406年）] \eventanchor{JH-0406-LIU-YU-DEFEATS-HUAN-REMNANTS}{东晋·刘裕等平定桓玄余部，东晋重新控制荆州与长江中游主要据点}
  \item[义熙三年（407年）] \eventanchor{JH-0407-NORTHERN-YAN}{北燕·冯跋推翻高云并在龙城建立北燕，后燕残余在东北改组}
  \item[义熙三年（407年）] \eventanchor{JH-0407-HELIAN-BOBO-XIA}{夏·赫连勃勃在朔方起兵建夏，利用铁弗部及边地军事网络挑战后秦}
  \item[义熙五年（409年）] \eventanchor{JH-0409-LIU-YU-SOUTHERN-YAN}{东晋与南燕·刘裕北伐南燕，东晋军跨淮河进攻广固}
  \item[义熙六年（410年）] \eventanchor{JH-0410-SOUTHERN-YAN-FALL}{东晋与南燕·刘裕攻陷广固，俘杀慕容超，南燕灭亡}
  \item[义熙八年（412年）] \eventanchor{JH-0412-LIU-YU-DEFEATS-LIU-YI}{东晋·刘裕击败并杀害据荆州的刘毅，将长江中游强镇纳入自身控制}
  \item[义熙九年（413年）] \eventanchor{JH-0413-QIAO-ZONG-DEFEATED}{东晋与谯蜀·朱龄石等东晋军攻灭谯蜀，成都重新纳入东晋治下}
  \item[义熙九年（413年）] \eventanchor{JH-0413-LU-XUN-SUPPRESSED}{东晋与岭南地方集团·卢循在广州一带败亡，孙恩—卢循集团的长期反乱结束}
  \item[义熙十年（414年）] \eventanchor{JH-0414-GWANGGAETO-STELE}{高句丽与东北诸政权·高句丽为好太王立碑，铭刻其征伐、城郭和王权记忆}
  \item[义熙十年（414年）] \eventanchor{JH-0414-WESTERN-QIN-DESTROYS-SOUTHERN-LIANG}{西秦与南凉·乞伏炽磐攻灭南凉，西秦在河湟地区扩展控制}
  \item[义熙十二年（416年）] \eventanchor{JH-0416-LIU-YU-NORTHERN-EXPEDITION}{东晋与后秦·刘裕自建康发动大规模北伐，水陆并进攻向后秦控制的河南与关中}
  \item[义熙十三年（417年）] \eventanchor{JH-0417-LIU-YU-TAKES-LUOYANG}{东晋与后秦·刘裕军攻取洛阳，后秦在河南的防线崩解}
  \item[义熙十三年（417年）] \eventanchor{JH-0417-LIU-YU-TAKES-CHANGAN}{东晋与后秦·东晋军攻入长安，姚泓降服，后秦灭亡}
  \item[义熙十四年（418年）] \eventanchor{JH-0418-XIA-TAKES-CHANGAN}{夏与东晋关中守军·赫连勃勃趁刘裕南返攻取长安，东晋失去关中据点}
  \item[建弘元年（420年）] \eventanchor{JH-0420-BINGLING-CAVE-169}{西秦·炳灵寺第169窟题记留下西秦年号与佛教造像活动的纪年证据}
  \item[泰常七年（422年）] \eventanchor{JH-0422-MINGYUAN-EMPEROR-DEATH}{北魏·明元帝拓跋嗣去世，太武帝拓跋焘继位，北魏军政权力完成交接}
  \item[始光元年（424年）] \eventanchor{JH-0424-TAIWU-CONSOLIDATES}{北魏·太武帝亲政后处置权臣与宗室竞争，重新集中北魏军政指挥}
  \item[始光四年（427年）] \eventanchor{JH-0427-NORTHERN-WEI-TAKES-TONGWAN}{北魏与夏·太武帝攻取夏国都城统万城，夏政权的核心基地失守}
  \item[元嘉七年（430年）] \eventanchor{JH-0430-LIU-SONG-NORTHERN-EXPEDITION}{刘宋与北魏·刘宋北伐一度收复黄河以南若干城镇，旋因北魏反击而难以久守}
  \item[神麚四年（431年）] \eventanchor{JH-0431-XIA-DESTROYED}{北魏与夏残余·赫连定败亡，夏国残余政权终结，北魏进一步控制关陇北部}
  \item[延和二年（433年）] \eventanchor{JH-0433-HELIAN-DING-CAPTURED}{北魏、夏残余与吐谷浑·吐谷浑将赫连定送交北魏，夏国最后君主的流亡政治结束}
  \item[太延二年（436年）] \eventanchor{JH-0436-NORTHERN-WEI-DESTROYS-NORTHERN-YAN}{北魏与北燕·北魏攻灭北燕，冯弘出奔高句丽，龙城政权结束}
  \item[太康二年（281年）] \eventanchor{JH-LEDGER-001-1}{西晋·“汲郡魏襄王墓出土竹书，朝廷命荀勖、和峤等整理其书写材料”的前期动员与资源准备}
  \item[太康二年（281年）] \eventanchor{JH-LEDGER-001-2}{西晋·“汲郡魏襄王墓出土竹书，朝廷命荀勖、和峤等整理其书写材料”的命令、联盟或地方响应}
  \item[太康二年（281年）] \eventanchor{JH-LEDGER-001-3}{西晋·汲郡魏襄王墓出土竹书，朝廷命荀勖、和峤等整理其书写材料}
  \item[太康二年（281年）] \eventanchor{JH-LEDGER-001-4}{西晋·“汲郡魏襄王墓出土竹书，朝廷命荀勖、和峤等整理其书写材料”的执行中的空间与人员处置}
  \item[太康二年（281年）] \eventanchor{JH-LEDGER-001-5}{西晋·“汲郡魏襄王墓出土竹书，朝廷命荀勖、和峤等整理其书写材料”的后续制度或军政调整}
  \item[太康二年（281年）] \eventanchor{JH-LEDGER-001-6}{西晋·“汲郡魏襄王墓出土竹书，朝廷命荀勖、和峤等整理其书写材料”的异源材料与影响范围的复核}
  \item[太熙元年（289年）] \eventanchor{JH-LEDGER-002-1}{西晋·“武帝病重时以外戚杨骏受遗诏辅政，围绕太子司马衷的继承安排集中于宫廷”的前期动员与资源准备}
  \item[太熙元年（289年）] \eventanchor{JH-LEDGER-002-2}{西晋·“武帝病重时以外戚杨骏受遗诏辅政，围绕太子司马衷的继承安排集中于宫廷”的命令、联盟或地方响应}
  \item[太熙元年（289年）] \eventanchor{JH-LEDGER-002-3}{西晋·武帝病重时以外戚杨骏受遗诏辅政，围绕太子司马衷的继承安排集中于宫廷}
  \item[太熙元年（289年）] \eventanchor{JH-LEDGER-002-4}{西晋·“武帝病重时以外戚杨骏受遗诏辅政，围绕太子司马衷的继承安排集中于宫廷”的执行中的空间与人员处置}
  \item[太熙元年（289年）] \eventanchor{JH-LEDGER-002-5}{西晋·“武帝病重时以外戚杨骏受遗诏辅政，围绕太子司马衷的继承安排集中于宫廷”的后续制度或军政调整}
  \item[太熙元年（289年）] \eventanchor{JH-LEDGER-002-6}{西晋·“武帝病重时以外戚杨骏受遗诏辅政，围绕太子司马衷的继承安排集中于宫廷”的异源材料与影响范围的复核}
  \item[太熙元年四月（290年）] \eventanchor{JH-LEDGER-003-1}{西晋·“司马炎去世，惠帝司马衷即位，杨骏以太傅录尚书事主持朝政”的前期动员与资源准备}
  \item[太熙元年四月（290年）] \eventanchor{JH-LEDGER-003-2}{西晋·“司马炎去世，惠帝司马衷即位，杨骏以太傅录尚书事主持朝政”的命令、联盟或地方响应}
  \item[太熙元年四月（290年）] \eventanchor{JH-LEDGER-003-3}{西晋·司马炎去世，惠帝司马衷即位，杨骏以太傅录尚书事主持朝政}
  \item[太熙元年四月（290年）] \eventanchor{JH-LEDGER-003-4}{西晋·“司马炎去世，惠帝司马衷即位，杨骏以太傅录尚书事主持朝政”的执行中的空间与人员处置}
  \item[太熙元年四月（290年）] \eventanchor{JH-LEDGER-003-5}{西晋·“司马炎去世，惠帝司马衷即位，杨骏以太傅录尚书事主持朝政”的后续制度或军政调整}
  \item[太熙元年四月（290年）] \eventanchor{JH-LEDGER-003-6}{西晋·“司马炎去世，惠帝司马衷即位，杨骏以太傅录尚书事主持朝政”的异源材料与影响范围的复核}
  \item[永平元年（291年）] \eventanchor{JH-LEDGER-004-1}{西晋·“贾后联络楚王司马玮等发动政变，杨骏被杀，杨氏外戚集团被清除”的前期动员与资源准备}
  \item[永平元年（291年）] \eventanchor{JH-LEDGER-004-2}{西晋·“贾后联络楚王司马玮等发动政变，杨骏被杀，杨氏外戚集团被清除”的命令、联盟或地方响应}
  \item[永平元年（291年）] \eventanchor{JH-LEDGER-004-3}{西晋·贾后联络楚王司马玮等发动政变，杨骏被杀，杨氏外戚集团被清除}
  \item[永平元年（291年）] \eventanchor{JH-LEDGER-004-4}{西晋·“贾后联络楚王司马玮等发动政变，杨骏被杀，杨氏外戚集团被清除”的执行中的空间与人员处置}
  \item[永平元年（291年）] \eventanchor{JH-LEDGER-004-5}{西晋·“贾后联络楚王司马玮等发动政变，杨骏被杀，杨氏外戚集团被清除”的后续制度或军政调整}
  \item[永平元年（291年）] \eventanchor{JH-LEDGER-004-6}{西晋·“贾后联络楚王司马玮等发动政变，杨骏被杀，杨氏外戚集团被清除”的异源材料与影响范围的复核}
  \item[元康元年（291年）] \eventanchor{JH-LEDGER-005-1}{西晋·“汝南王司马亮与卫瓘辅政后又被贾后借楚王司马玮诛杀，司马玮旋即被处死”的前期动员与资源准备}
  \item[元康元年（291年）] \eventanchor{JH-LEDGER-005-2}{西晋·“汝南王司马亮与卫瓘辅政后又被贾后借楚王司马玮诛杀，司马玮旋即被处死”的命令、联盟或地方响应}
  \item[元康元年（291年）] \eventanchor{JH-LEDGER-005-3}{西晋·汝南王司马亮与卫瓘辅政后又被贾后借楚王司马玮诛杀，司马玮旋即被处死}
  \item[元康元年（291年）] \eventanchor{JH-LEDGER-005-4}{西晋·“汝南王司马亮与卫瓘辅政后又被贾后借楚王司马玮诛杀，司马玮旋即被处死”的执行中的空间与人员处置}
  \item[元康元年（291年）] \eventanchor{JH-LEDGER-005-5}{西晋·“汝南王司马亮与卫瓘辅政后又被贾后借楚王司马玮诛杀，司马玮旋即被处死”的后续制度或军政调整}
  \item[元康元年（291年）] \eventanchor{JH-LEDGER-005-6}{西晋·“汝南王司马亮与卫瓘辅政后又被贾后借楚王司马玮诛杀，司马玮旋即被处死”的异源材料与影响范围的复核}
  \item[元康六年至九年（296年至299年）] \eventanchor{JH-LEDGER-006-1}{西晋与关陇氐羌集团·“氐族齐万年在关中起兵，西晋多次征讨并长期调动关陇军粮，至299年叛乱被平定”的前期动员与资源准备}
  \item[元康六年至九年（296年至299年）] \eventanchor{JH-LEDGER-006-2}{西晋与关陇氐羌集团·“氐族齐万年在关中起兵，西晋多次征讨并长期调动关陇军粮，至299年叛乱被平定”的命令、联盟或地方响应}
  \item[元康六年至九年（296年至299年）] \eventanchor{JH-LEDGER-006-3}{西晋与关陇氐羌集团·氐族齐万年在关中起兵，西晋多次征讨并长期调动关陇军粮，至299年叛乱被平定}
  \item[元康六年至九年（296年至299年）] \eventanchor{JH-LEDGER-006-4}{西晋与关陇氐羌集团·“氐族齐万年在关中起兵，西晋多次征讨并长期调动关陇军粮，至299年叛乱被平定”的执行中的空间与人员处置}
  \item[元康六年至九年（296年至299年）] \eventanchor{JH-LEDGER-006-5}{西晋与关陇氐羌集团·“氐族齐万年在关中起兵，西晋多次征讨并长期调动关陇军粮，至299年叛乱被平定”的后续制度或军政调整}
  \item[元康六年至九年（296年至299年）] \eventanchor{JH-LEDGER-006-6}{西晋与关陇氐羌集团·“氐族齐万年在关中起兵，西晋多次征讨并长期调动关陇军粮，至299年叛乱被平定”的异源材料与影响范围的复核}
  \item[永康元年（300年）] \eventanchor{JH-LEDGER-007-1}{西晋·“贾后以谋反罪废杀太子司马遹，皇位继承与宫廷联盟由此公开破裂”的前期动员与资源准备}
  \item[永康元年（300年）] \eventanchor{JH-LEDGER-007-2}{西晋·“贾后以谋反罪废杀太子司马遹，皇位继承与宫廷联盟由此公开破裂”的命令、联盟或地方响应}
  \item[永康元年（300年）] \eventanchor{JH-LEDGER-007-3}{西晋·贾后以谋反罪废杀太子司马遹，皇位继承与宫廷联盟由此公开破裂}
  \item[永康元年（300年）] \eventanchor{JH-LEDGER-007-4}{西晋·“贾后以谋反罪废杀太子司马遹，皇位继承与宫廷联盟由此公开破裂”的执行中的空间与人员处置}
  \item[永康元年（300年）] \eventanchor{JH-LEDGER-007-5}{西晋·“贾后以谋反罪废杀太子司马遹，皇位继承与宫廷联盟由此公开破裂”的后续制度或军政调整}
  \item[永康元年（300年）] \eventanchor{JH-LEDGER-007-6}{西晋·“贾后以谋反罪废杀太子司马遹，皇位继承与宫廷联盟由此公开破裂”的异源材料与影响范围的复核}
  \item[永宁元年（301年）] \eventanchor{JH-LEDGER-008-1}{西晋·“赵王司马伦杀贾后并废惠帝自立为帝，短暂改变西晋帝统”的前期动员与资源准备}
  \item[永宁元年（301年）] \eventanchor{JH-LEDGER-008-2}{西晋·“赵王司马伦杀贾后并废惠帝自立为帝，短暂改变西晋帝统”的命令、联盟或地方响应}
  \item[永宁元年（301年）] \eventanchor{JH-LEDGER-008-3}{西晋·赵王司马伦杀贾后并废惠帝自立为帝，短暂改变西晋帝统}
  \item[永宁元年（301年）] \eventanchor{JH-LEDGER-008-4}{西晋·“赵王司马伦杀贾后并废惠帝自立为帝，短暂改变西晋帝统”的执行中的空间与人员处置}
  \item[永宁元年（301年）] \eventanchor{JH-LEDGER-008-5}{西晋·“赵王司马伦杀贾后并废惠帝自立为帝，短暂改变西晋帝统”的后续制度或军政调整}
  \item[永宁元年（301年）] \eventanchor{JH-LEDGER-008-6}{西晋·“赵王司马伦杀贾后并废惠帝自立为帝，短暂改变西晋帝统”的异源材料与影响范围的复核}
  \item[永宁元年（301年）] \eventanchor{JH-LEDGER-009-1}{西晋·“齐王司马冏等起兵推翻司马伦，复立惠帝并以大司马辅政”的前期动员与资源准备}
  \item[永宁元年（301年）] \eventanchor{JH-LEDGER-009-2}{西晋·“齐王司马冏等起兵推翻司马伦，复立惠帝并以大司马辅政”的命令、联盟或地方响应}
  \item[永宁元年（301年）] \eventanchor{JH-LEDGER-009-3}{西晋·齐王司马冏等起兵推翻司马伦，复立惠帝并以大司马辅政}
  \item[永宁元年（301年）] \eventanchor{JH-LEDGER-009-4}{西晋·“齐王司马冏等起兵推翻司马伦，复立惠帝并以大司马辅政”的执行中的空间与人员处置}
  \item[永宁元年（301年）] \eventanchor{JH-LEDGER-009-5}{西晋·“齐王司马冏等起兵推翻司马伦，复立惠帝并以大司马辅政”的后续制度或军政调整}
  \item[永宁元年（301年）] \eventanchor{JH-LEDGER-009-6}{西晋·“齐王司马冏等起兵推翻司马伦，复立惠帝并以大司马辅政”的异源材料与影响范围的复核}
  \item[太安元年（302年）] \eventanchor{JH-LEDGER-010-1}{西晋·“河间王司马颙与成都王司马颖起兵，长沙王司马乂在洛阳杀齐王司马冏并接掌朝政”的前期动员与资源准备}
  \item[太安元年（302年）] \eventanchor{JH-LEDGER-010-2}{西晋·“河间王司马颙与成都王司马颖起兵，长沙王司马乂在洛阳杀齐王司马冏并接掌朝政”的命令、联盟或地方响应}
  \item[太安元年（302年）] \eventanchor{JH-LEDGER-010-3}{西晋·河间王司马颙与成都王司马颖起兵，长沙王司马乂在洛阳杀齐王司马冏并接掌朝政}
  \item[太安元年（302年）] \eventanchor{JH-LEDGER-010-4}{西晋·“河间王司马颙与成都王司马颖起兵，长沙王司马乂在洛阳杀齐王司马冏并接掌朝政”的执行中的空间与人员处置}
  \item[太安元年（302年）] \eventanchor{JH-LEDGER-010-5}{西晋·“河间王司马颙与成都王司马颖起兵，长沙王司马乂在洛阳杀齐王司马冏并接掌朝政”的后续制度或军政调整}
  \item[太安元年（302年）] \eventanchor{JH-LEDGER-010-6}{西晋·“河间王司马颙与成都王司马颖起兵，长沙王司马乂在洛阳杀齐王司马冏并接掌朝政”的异源材料与影响范围的复核}
  \item[太安二年（303年）] \eventanchor{JH-LEDGER-011-1}{西晋·“长沙王司马乂据洛阳抵抗河间王司马颙和成都王司马颖联军，京师陷入长期战斗”的前期动员与资源准备}
  \item[太安二年（303年）] \eventanchor{JH-LEDGER-011-2}{西晋·“长沙王司马乂据洛阳抵抗河间王司马颙和成都王司马颖联军，京师陷入长期战斗”的命令、联盟或地方响应}
  \item[太安二年（303年）] \eventanchor{JH-LEDGER-011-3}{西晋·长沙王司马乂据洛阳抵抗河间王司马颙和成都王司马颖联军，京师陷入长期战斗}
  \item[太安二年（303年）] \eventanchor{JH-LEDGER-011-4}{西晋·“长沙王司马乂据洛阳抵抗河间王司马颙和成都王司马颖联军，京师陷入长期战斗”的执行中的空间与人员处置}
  \item[太安二年（303年）] \eventanchor{JH-LEDGER-011-5}{西晋·“长沙王司马乂据洛阳抵抗河间王司马颙和成都王司马颖联军，京师陷入长期战斗”的后续制度或军政调整}
  \item[太安二年（303年）] \eventanchor{JH-LEDGER-011-6}{西晋·“长沙王司马乂据洛阳抵抗河间王司马颙和成都王司马颖联军，京师陷入长期战斗”的异源材料与影响范围的复核}
  \item[太安二年（303年）] \eventanchor{JH-LEDGER-012-1}{成汉与益州西晋守军·“巴氐首领李雄攻取成都，益州的西晋统治被成汉政权取代”的前期动员与资源准备}
  \item[太安二年（303年）] \eventanchor{JH-LEDGER-012-2}{成汉与益州西晋守军·“巴氐首领李雄攻取成都，益州的西晋统治被成汉政权取代”的命令、联盟或地方响应}
  \item[太安二年（303年）] \eventanchor{JH-LEDGER-012-3}{成汉与益州西晋守军·巴氐首领李雄攻取成都，益州的西晋统治被成汉政权取代}
  \item[太安二年（303年）] \eventanchor{JH-LEDGER-012-4}{成汉与益州西晋守军·“巴氐首领李雄攻取成都，益州的西晋统治被成汉政权取代”的执行中的空间与人员处置}
  \item[太安二年（303年）] \eventanchor{JH-LEDGER-012-5}{成汉与益州西晋守军·“巴氐首领李雄攻取成都，益州的西晋统治被成汉政权取代”的后续制度或军政调整}
  \item[太安二年（303年）] \eventanchor{JH-LEDGER-012-6}{成汉与益州西晋守军·“巴氐首领李雄攻取成都，益州的西晋统治被成汉政权取代”的异源材料与影响范围的复核}
  \item[永安元年（304年）] \eventanchor{JH-LEDGER-013-1}{西晋·“司马乂被东海王司马越等交给河间王司马颙军并遭杀害，洛阳政权落入外军支配”的前期动员与资源准备}
  \item[永安元年（304年）] \eventanchor{JH-LEDGER-013-2}{西晋·“司马乂被东海王司马越等交给河间王司马颙军并遭杀害，洛阳政权落入外军支配”的命令、联盟或地方响应}
  \item[永安元年（304年）] \eventanchor{JH-LEDGER-013-3}{西晋·司马乂被东海王司马越等交给河间王司马颙军并遭杀害，洛阳政权落入外军支配}
  \item[永安元年（304年）] \eventanchor{JH-LEDGER-013-4}{西晋·“司马乂被东海王司马越等交给河间王司马颙军并遭杀害，洛阳政权落入外军支配”的执行中的空间与人员处置}
  \item[永安元年（304年）] \eventanchor{JH-LEDGER-013-5}{西晋·“司马乂被东海王司马越等交给河间王司马颙军并遭杀害，洛阳政权落入外军支配”的后续制度或军政调整}
  \item[永安元年（304年）] \eventanchor{JH-LEDGER-013-6}{西晋·“司马乂被东海王司马越等交给河间王司马颙军并遭杀害，洛阳政权落入外军支配”的异源材料与影响范围的复核}
  \item[永兴元年（304年）] \eventanchor{JH-LEDGER-014-1}{汉赵与并州诸部·“刘渊在左国城起兵，自称汉王，汇聚匈奴五部及并州反晋力量”的前期动员与资源准备}
  \item[永兴元年（304年）] \eventanchor{JH-LEDGER-014-2}{汉赵与并州诸部·“刘渊在左国城起兵，自称汉王，汇聚匈奴五部及并州反晋力量”的命令、联盟或地方响应}
  \item[永兴元年（304年）] \eventanchor{JH-LEDGER-014-3}{汉赵与并州诸部·刘渊在左国城起兵，自称汉王，汇聚匈奴五部及并州反晋力量}
  \item[永兴元年（304年）] \eventanchor{JH-LEDGER-014-4}{汉赵与并州诸部·“刘渊在左国城起兵，自称汉王，汇聚匈奴五部及并州反晋力量”的执行中的空间与人员处置}
  \item[永兴元年（304年）] \eventanchor{JH-LEDGER-014-5}{汉赵与并州诸部·“刘渊在左国城起兵，自称汉王，汇聚匈奴五部及并州反晋力量”的后续制度或军政调整}
  \item[永兴元年（304年）] \eventanchor{JH-LEDGER-014-6}{汉赵与并州诸部·“刘渊在左国城起兵，自称汉王，汇聚匈奴五部及并州反晋力量”的异源材料与影响范围的复核}
  \item[永兴元年（304年）] \eventanchor{JH-LEDGER-015-1}{西晋与江东地方集团·“陈敏据江东起兵并自置官属，试图利用八王之乱切割东南地区”的前期动员与资源准备}
  \item[永兴元年（304年）] \eventanchor{JH-LEDGER-015-2}{西晋与江东地方集团·“陈敏据江东起兵并自置官属，试图利用八王之乱切割东南地区”的命令、联盟或地方响应}
  \item[永兴元年（304年）] \eventanchor{JH-LEDGER-015-3}{西晋与江东地方集团·陈敏据江东起兵并自置官属，试图利用八王之乱切割东南地区}
  \item[永兴元年（304年）] \eventanchor{JH-LEDGER-015-4}{西晋与江东地方集团·“陈敏据江东起兵并自置官属，试图利用八王之乱切割东南地区”的执行中的空间与人员处置}
  \item[永兴元年（304年）] \eventanchor{JH-LEDGER-015-5}{西晋与江东地方集团·“陈敏据江东起兵并自置官属，试图利用八王之乱切割东南地区”的后续制度或军政调整}
  \item[永兴元年（304年）] \eventanchor{JH-LEDGER-015-6}{西晋与江东地方集团·“陈敏据江东起兵并自置官属，试图利用八王之乱切割东南地区”的异源材料与影响范围的复核}
  \item[永兴二年（305年）] \eventanchor{JH-LEDGER-016-1}{西晋·“东海王司马越以讨河间王司马颙为名起兵，八王之乱扩大为跨州郡的长期战争”的前期动员与资源准备}
  \item[永兴二年（305年）] \eventanchor{JH-LEDGER-016-2}{西晋·“东海王司马越以讨河间王司马颙为名起兵，八王之乱扩大为跨州郡的长期战争”的命令、联盟或地方响应}
  \item[永兴二年（305年）] \eventanchor{JH-LEDGER-016-3}{西晋·东海王司马越以讨河间王司马颙为名起兵，八王之乱扩大为跨州郡的长期战争}
  \item[永兴二年（305年）] \eventanchor{JH-LEDGER-016-4}{西晋·“东海王司马越以讨河间王司马颙为名起兵，八王之乱扩大为跨州郡的长期战争”的执行中的空间与人员处置}
  \item[永兴二年（305年）] \eventanchor{JH-LEDGER-016-5}{西晋·“东海王司马越以讨河间王司马颙为名起兵，八王之乱扩大为跨州郡的长期战争”的后续制度或军政调整}
  \item[永兴二年（305年）] \eventanchor{JH-LEDGER-016-6}{西晋·“东海王司马越以讨河间王司马颙为名起兵，八王之乱扩大为跨州郡的长期战争”的异源材料与影响范围的复核}
  \item[光熙元年（306年）] \eventanchor{JH-LEDGER-017-1}{西晋·“惠帝在洛阳去世，司马炽即位为怀帝，司马越继续掌握主要军政权”的前期动员与资源准备}
  \item[光熙元年（306年）] \eventanchor{JH-LEDGER-017-2}{西晋·“惠帝在洛阳去世，司马炽即位为怀帝，司马越继续掌握主要军政权”的命令、联盟或地方响应}
  \item[光熙元年（306年）] \eventanchor{JH-LEDGER-017-3}{西晋·惠帝在洛阳去世，司马炽即位为怀帝，司马越继续掌握主要军政权}
  \item[光熙元年（306年）] \eventanchor{JH-LEDGER-017-4}{西晋·“惠帝在洛阳去世，司马炽即位为怀帝，司马越继续掌握主要军政权”的执行中的空间与人员处置}
  \item[光熙元年（306年）] \eventanchor{JH-LEDGER-017-5}{西晋·“惠帝在洛阳去世，司马炽即位为怀帝，司马越继续掌握主要军政权”的后续制度或军政调整}
  \item[光熙元年（306年）] \eventanchor{JH-LEDGER-017-6}{西晋·“惠帝在洛阳去世，司马炽即位为怀帝，司马越继续掌握主要军政权”的异源材料与影响范围的复核}
  \item[晏平元年（306年）] \eventanchor{JH-LEDGER-018-1}{成汉·“李雄在成都称帝，国号成，使巴蜀割据由起事集团转为帝国体制”的前期动员与资源准备}
  \item[晏平元年（306年）] \eventanchor{JH-LEDGER-018-2}{成汉·“李雄在成都称帝，国号成，使巴蜀割据由起事集团转为帝国体制”的命令、联盟或地方响应}
  \item[晏平元年（306年）] \eventanchor{JH-LEDGER-018-3}{成汉·李雄在成都称帝，国号成，使巴蜀割据由起事集团转为帝国体制}
  \item[晏平元年（306年）] \eventanchor{JH-LEDGER-018-4}{成汉·“李雄在成都称帝，国号成，使巴蜀割据由起事集团转为帝国体制”的执行中的空间与人员处置}
  \item[晏平元年（306年）] \eventanchor{JH-LEDGER-018-5}{成汉·“李雄在成都称帝，国号成，使巴蜀割据由起事集团转为帝国体制”的后续制度或军政调整}
  \item[晏平元年（306年）] \eventanchor{JH-LEDGER-018-6}{成汉·“李雄在成都称帝，国号成，使巴蜀割据由起事集团转为帝国体制”的异源材料与影响范围的复核}
  \item[永嘉元年（307年）] \eventanchor{JH-LEDGER-019-1}{西晋江东·“司马睿出镇建业，王导等协助联络江东士族并重整东南军政”的前期动员与资源准备}
  \item[永嘉元年（307年）] \eventanchor{JH-LEDGER-019-2}{西晋江东·“司马睿出镇建业，王导等协助联络江东士族并重整东南军政”的命令、联盟或地方响应}
  \item[永嘉元年（307年）] \eventanchor{JH-LEDGER-019-3}{西晋江东·司马睿出镇建业，王导等协助联络江东士族并重整东南军政}
  \item[永嘉元年（307年）] \eventanchor{JH-LEDGER-019-4}{西晋江东·“司马睿出镇建业，王导等协助联络江东士族并重整东南军政”的执行中的空间与人员处置}
  \item[永嘉元年（307年）] \eventanchor{JH-LEDGER-019-5}{西晋江东·“司马睿出镇建业，王导等协助联络江东士族并重整东南军政”的后续制度或军政调整}
  \item[永嘉元年（307年）] \eventanchor{JH-LEDGER-019-6}{西晋江东·“司马睿出镇建业，王导等协助联络江东士族并重整东南军政”的异源材料与影响范围的复核}
  \item[永凤元年（308年）] \eventanchor{JH-LEDGER-020-1}{汉赵·“刘渊称帝，正式以汉为国号并以平阳为政治中心”的前期动员与资源准备}
  \item[永凤元年（308年）] \eventanchor{JH-LEDGER-020-2}{汉赵·“刘渊称帝，正式以汉为国号并以平阳为政治中心”的命令、联盟或地方响应}
  \item[永凤元年（308年）] \eventanchor{JH-LEDGER-020-3}{汉赵·刘渊称帝，正式以汉为国号并以平阳为政治中心}
  \item[永凤元年（308年）] \eventanchor{JH-LEDGER-020-4}{汉赵·“刘渊称帝，正式以汉为国号并以平阳为政治中心”的执行中的空间与人员处置}
  \item[永凤元年（308年）] \eventanchor{JH-LEDGER-020-5}{汉赵·“刘渊称帝，正式以汉为国号并以平阳为政治中心”的后续制度或军政调整}
  \item[永凤元年（308年）] \eventanchor{JH-LEDGER-020-6}{汉赵·“刘渊称帝，正式以汉为国号并以平阳为政治中心”的异源材料与影响范围的复核}
  \item[永嘉三年（309年）] \eventanchor{JH-LEDGER-021-1}{汉赵与西晋·“刘聪、王弥等围攻洛阳，西晋京师依靠临时征兵和城防勉强守住”的前期动员与资源准备}
  \item[永嘉三年（309年）] \eventanchor{JH-LEDGER-021-2}{汉赵与西晋·“刘聪、王弥等围攻洛阳，西晋京师依靠临时征兵和城防勉强守住”的命令、联盟或地方响应}
  \item[永嘉三年（309年）] \eventanchor{JH-LEDGER-021-3}{汉赵与西晋·刘聪、王弥等围攻洛阳，西晋京师依靠临时征兵和城防勉强守住}
  \item[永嘉三年（309年）] \eventanchor{JH-LEDGER-021-4}{汉赵与西晋·“刘聪、王弥等围攻洛阳，西晋京师依靠临时征兵和城防勉强守住”的执行中的空间与人员处置}
  \item[永嘉三年（309年）] \eventanchor{JH-LEDGER-021-5}{汉赵与西晋·“刘聪、王弥等围攻洛阳，西晋京师依靠临时征兵和城防勉强守住”的后续制度或军政调整}
  \item[永嘉三年（309年）] \eventanchor{JH-LEDGER-021-6}{汉赵与西晋·“刘聪、王弥等围攻洛阳，西晋京师依靠临时征兵和城防勉强守住”的异源材料与影响范围的复核}
  \item[河瑞二年（310年）] \eventanchor{JH-LEDGER-022-1}{汉赵·“刘渊去世后刘聪经宫廷斗争继位，汉赵继续对西晋发动攻势”的前期动员与资源准备}
  \item[河瑞二年（310年）] \eventanchor{JH-LEDGER-022-2}{汉赵·“刘渊去世后刘聪经宫廷斗争继位，汉赵继续对西晋发动攻势”的命令、联盟或地方响应}
  \item[河瑞二年（310年）] \eventanchor{JH-LEDGER-022-3}{汉赵·刘渊去世后刘聪经宫廷斗争继位，汉赵继续对西晋发动攻势}
  \item[河瑞二年（310年）] \eventanchor{JH-LEDGER-022-4}{汉赵·“刘渊去世后刘聪经宫廷斗争继位，汉赵继续对西晋发动攻势”的执行中的空间与人员处置}
  \item[河瑞二年（310年）] \eventanchor{JH-LEDGER-022-5}{汉赵·“刘渊去世后刘聪经宫廷斗争继位，汉赵继续对西晋发动攻势”的后续制度或军政调整}
  \item[河瑞二年（310年）] \eventanchor{JH-LEDGER-022-6}{汉赵·“刘渊去世后刘聪经宫廷斗争继位，汉赵继续对西晋发动攻势”的异源材料与影响范围的复核}
  \item[永嘉四年（310年）] \eventanchor{JH-LEDGER-023-1}{西晋·“司马越在东出避乱途中去世，其部众和辎重随即失去统一指挥”的前期动员与资源准备}
  \item[永嘉四年（310年）] \eventanchor{JH-LEDGER-023-2}{西晋·“司马越在东出避乱途中去世，其部众和辎重随即失去统一指挥”的命令、联盟或地方响应}
  \item[永嘉四年（310年）] \eventanchor{JH-LEDGER-023-3}{西晋·司马越在东出避乱途中去世，其部众和辎重随即失去统一指挥}
  \item[永嘉四年（310年）] \eventanchor{JH-LEDGER-023-4}{西晋·“司马越在东出避乱途中去世，其部众和辎重随即失去统一指挥”的执行中的空间与人员处置}
  \item[永嘉四年（310年）] \eventanchor{JH-LEDGER-023-5}{西晋·“司马越在东出避乱途中去世，其部众和辎重随即失去统一指挥”的后续制度或军政调整}
  \item[永嘉四年（310年）] \eventanchor{JH-LEDGER-023-6}{西晋·“司马越在东出避乱途中去世，其部众和辎重随即失去统一指挥”的异源材料与影响范围的复核}
  \item[永嘉五年六月（311年）] \eventanchor{JH-LEDGER-024-1}{汉赵与西晋·“汉赵军攻陷洛阳，大量官民死散，西晋都城体系遭到毁灭性破坏”的前期动员与资源准备}
  \item[永嘉五年六月（311年）] \eventanchor{JH-LEDGER-024-2}{汉赵与西晋·“汉赵军攻陷洛阳，大量官民死散，西晋都城体系遭到毁灭性破坏”的命令、联盟或地方响应}
  \item[永嘉五年六月（311年）] \eventanchor{JH-LEDGER-024-3}{汉赵与西晋·汉赵军攻陷洛阳，大量官民死散，西晋都城体系遭到毁灭性破坏}
  \item[永嘉五年六月（311年）] \eventanchor{JH-LEDGER-024-4}{汉赵与西晋·“汉赵军攻陷洛阳，大量官民死散，西晋都城体系遭到毁灭性破坏”的执行中的空间与人员处置}
  \item[永嘉五年六月（311年）] \eventanchor{JH-LEDGER-024-5}{汉赵与西晋·“汉赵军攻陷洛阳，大量官民死散，西晋都城体系遭到毁灭性破坏”的后续制度或军政调整}
  \item[永嘉五年六月（311年）] \eventanchor{JH-LEDGER-024-6}{汉赵与西晋·“汉赵军攻陷洛阳，大量官民死散，西晋都城体系遭到毁灭性破坏”的异源材料与影响范围的复核}
  \item[永嘉五年（311年）] \eventanchor{JH-LEDGER-025-1}{汉赵与西晋·“怀帝在洛阳陷落时被汉赵俘获并押往平阳，西晋皇统陷入屈辱性危机”的前期动员与资源准备}
  \item[永嘉五年（311年）] \eventanchor{JH-LEDGER-025-2}{汉赵与西晋·“怀帝在洛阳陷落时被汉赵俘获并押往平阳，西晋皇统陷入屈辱性危机”的命令、联盟或地方响应}
  \item[永嘉五年（311年）] \eventanchor{JH-LEDGER-025-3}{汉赵与西晋·怀帝在洛阳陷落时被汉赵俘获并押往平阳，西晋皇统陷入屈辱性危机}
  \item[永嘉五年（311年）] \eventanchor{JH-LEDGER-025-4}{汉赵与西晋·“怀帝在洛阳陷落时被汉赵俘获并押往平阳，西晋皇统陷入屈辱性危机”的执行中的空间与人员处置}
  \item[永嘉五年（311年）] \eventanchor{JH-LEDGER-025-5}{汉赵与西晋·“怀帝在洛阳陷落时被汉赵俘获并押往平阳，西晋皇统陷入屈辱性危机”的后续制度或军政调整}
  \item[永嘉五年（311年）] \eventanchor{JH-LEDGER-025-6}{汉赵与西晋·“怀帝在洛阳陷落时被汉赵俘获并押往平阳，西晋皇统陷入屈辱性危机”的异源材料与影响范围的复核}
  \item[建兴元年（313年）] \eventanchor{JH-LEDGER-026-1}{汉赵与西晋·“汉赵杀害被俘的怀帝，长安的司马邺被拥立为愍帝”的前期动员与资源准备}
  \item[建兴元年（313年）] \eventanchor{JH-LEDGER-026-2}{汉赵与西晋·“汉赵杀害被俘的怀帝，长安的司马邺被拥立为愍帝”的命令、联盟或地方响应}
  \item[建兴元年（313年）] \eventanchor{JH-LEDGER-026-3}{汉赵与西晋·汉赵杀害被俘的怀帝，长安的司马邺被拥立为愍帝}
  \item[建兴元年（313年）] \eventanchor{JH-LEDGER-026-4}{汉赵与西晋·“汉赵杀害被俘的怀帝，长安的司马邺被拥立为愍帝”的执行中的空间与人员处置}
  \item[建兴元年（313年）] \eventanchor{JH-LEDGER-026-5}{汉赵与西晋·“汉赵杀害被俘的怀帝，长安的司马邺被拥立为愍帝”的后续制度或军政调整}
  \item[建兴元年（313年）] \eventanchor{JH-LEDGER-026-6}{汉赵与西晋·“汉赵杀害被俘的怀帝，长安的司马邺被拥立为愍帝”的异源材料与影响范围的复核}
  \item[建兴元年（313年）] \eventanchor{JH-LEDGER-027-1}{东晋前身与河南地方集团·“祖逖率部北渡长江，经屯田与结盟经营豫州、河南一带”的前期动员与资源准备}
  \item[建兴元年（313年）] \eventanchor{JH-LEDGER-027-2}{东晋前身与河南地方集团·“祖逖率部北渡长江，经屯田与结盟经营豫州、河南一带”的命令、联盟或地方响应}
  \item[建兴元年（313年）] \eventanchor{JH-LEDGER-027-3}{东晋前身与河南地方集团·祖逖率部北渡长江，经屯田与结盟经营豫州、河南一带}
  \item[建兴元年（313年）] \eventanchor{JH-LEDGER-027-4}{东晋前身与河南地方集团·“祖逖率部北渡长江，经屯田与结盟经营豫州、河南一带”的执行中的空间与人员处置}
  \item[建兴元年（313年）] \eventanchor{JH-LEDGER-027-5}{东晋前身与河南地方集团·“祖逖率部北渡长江，经屯田与结盟经营豫州、河南一带”的后续制度或军政调整}
  \item[建兴元年（313年）] \eventanchor{JH-LEDGER-027-6}{东晋前身与河南地方集团·“祖逖率部北渡长江，经屯田与结盟经营豫州、河南一带”的异源材料与影响范围的复核}
  \item[建兴三年（315年）] \eventanchor{JH-LEDGER-028-1}{代国与西晋残余政权·“拓跋猗卢受晋封为代王，在北部边地扩展部众与城郭控制”的前期动员与资源准备}
  \item[建兴三年（315年）] \eventanchor{JH-LEDGER-028-2}{代国与西晋残余政权·“拓跋猗卢受晋封为代王，在北部边地扩展部众与城郭控制”的命令、联盟或地方响应}
  \item[建兴三年（315年）] \eventanchor{JH-LEDGER-028-3}{代国与西晋残余政权·拓跋猗卢受晋封为代王，在北部边地扩展部众与城郭控制}
  \item[建兴三年（315年）] \eventanchor{JH-LEDGER-028-4}{代国与西晋残余政权·“拓跋猗卢受晋封为代王，在北部边地扩展部众与城郭控制”的执行中的空间与人员处置}
  \item[建兴三年（315年）] \eventanchor{JH-LEDGER-028-5}{代国与西晋残余政权·“拓跋猗卢受晋封为代王，在北部边地扩展部众与城郭控制”的后续制度或军政调整}
  \item[建兴三年（315年）] \eventanchor{JH-LEDGER-028-6}{代国与西晋残余政权·“拓跋猗卢受晋封为代王，在北部边地扩展部众与城郭控制”的异源材料与影响范围的复核}
  \item[建兴四年十一月（316年）] \eventanchor{JH-LEDGER-029-1}{汉赵与西晋·“刘曜攻陷长安，愍帝出降，西晋在北方的最后朝廷覆灭”的前期动员与资源准备}
  \item[建兴四年十一月（316年）] \eventanchor{JH-LEDGER-029-2}{汉赵与西晋·“刘曜攻陷长安，愍帝出降，西晋在北方的最后朝廷覆灭”的命令、联盟或地方响应}
  \item[建兴四年十一月（316年）] \eventanchor{JH-LEDGER-029-3}{汉赵与西晋·刘曜攻陷长安，愍帝出降，西晋在北方的最后朝廷覆灭}
  \item[建兴四年十一月（316年）] \eventanchor{JH-LEDGER-029-4}{汉赵与西晋·“刘曜攻陷长安，愍帝出降，西晋在北方的最后朝廷覆灭”的执行中的空间与人员处置}
  \item[建兴四年十一月（316年）] \eventanchor{JH-LEDGER-029-5}{汉赵与西晋·“刘曜攻陷长安，愍帝出降，西晋在北方的最后朝廷覆灭”的后续制度或军政调整}
  \item[建兴四年十一月（316年）] \eventanchor{JH-LEDGER-029-6}{汉赵与西晋·“刘曜攻陷长安，愍帝出降，西晋在北方的最后朝廷覆灭”的异源材料与影响范围的复核}
  \item[建武元年三月（317年）] \eventanchor{JH-LEDGER-030-1}{东晋·“司马睿在建康即晋王位，建立承继西晋的南方朝廷框架”的前期动员与资源准备}
  \item[建武元年三月（317年）] \eventanchor{JH-LEDGER-030-2}{东晋·“司马睿在建康即晋王位，建立承继西晋的南方朝廷框架”的命令、联盟或地方响应}
  \item[建武元年三月（317年）] \eventanchor{JH-LEDGER-030-3}{东晋·司马睿在建康即晋王位，建立承继西晋的南方朝廷框架}
  \item[建武元年三月（317年）] \eventanchor{JH-LEDGER-030-4}{东晋·“司马睿在建康即晋王位，建立承继西晋的南方朝廷框架”的执行中的空间与人员处置}
  \item[建武元年三月（317年）] \eventanchor{JH-LEDGER-030-5}{东晋·“司马睿在建康即晋王位，建立承继西晋的南方朝廷框架”的后续制度或军政调整}
  \item[建武元年三月（317年）] \eventanchor{JH-LEDGER-030-6}{东晋·“司马睿在建康即晋王位，建立承继西晋的南方朝廷框架”的异源材料与影响范围的复核}
  \item[太兴元年三月（318年）] \eventanchor{JH-LEDGER-031-1}{东晋·“司马睿称帝，是为元帝，东晋以建康为都正式确立”的前期动员与资源准备}
  \item[太兴元年三月（318年）] \eventanchor{JH-LEDGER-031-2}{东晋·“司马睿称帝，是为元帝，东晋以建康为都正式确立”的命令、联盟或地方响应}
  \item[太兴元年三月（318年）] \eventanchor{JH-LEDGER-031-3}{东晋·司马睿称帝，是为元帝，东晋以建康为都正式确立}
  \item[太兴元年三月（318年）] \eventanchor{JH-LEDGER-031-4}{东晋·“司马睿称帝，是为元帝，东晋以建康为都正式确立”的执行中的空间与人员处置}
  \item[太兴元年三月（318年）] \eventanchor{JH-LEDGER-031-5}{东晋·“司马睿称帝，是为元帝，东晋以建康为都正式确立”的后续制度或军政调整}
  \item[太兴元年三月（318年）] \eventanchor{JH-LEDGER-031-6}{东晋·“司马睿称帝，是为元帝，东晋以建康为都正式确立”的异源材料与影响范围的复核}
  \item[太兴二年（319年）] \eventanchor{JH-LEDGER-032-1}{后赵·“石勒在襄国称赵王并建立独立政权，后赵由汉赵将领集团分出”的前期动员与资源准备}
  \item[太兴二年（319年）] \eventanchor{JH-LEDGER-032-2}{后赵·“石勒在襄国称赵王并建立独立政权，后赵由汉赵将领集团分出”的命令、联盟或地方响应}
  \item[太兴二年（319年）] \eventanchor{JH-LEDGER-032-3}{后赵·石勒在襄国称赵王并建立独立政权，后赵由汉赵将领集团分出}
  \item[太兴二年（319年）] \eventanchor{JH-LEDGER-032-4}{后赵·“石勒在襄国称赵王并建立独立政权，后赵由汉赵将领集团分出”的执行中的空间与人员处置}
  \item[太兴二年（319年）] \eventanchor{JH-LEDGER-032-5}{后赵·“石勒在襄国称赵王并建立独立政权，后赵由汉赵将领集团分出”的后续制度或军政调整}
  \item[太兴二年（319年）] \eventanchor{JH-LEDGER-032-6}{后赵·“石勒在襄国称赵王并建立独立政权，后赵由汉赵将领集团分出”的异源材料与影响范围的复核}
  \item[光初二年（319年）] \eventanchor{JH-LEDGER-033-1}{前赵·“刘曜在长安改汉国号为赵，以关中为中心重建政权身份”的前期动员与资源准备}
  \item[光初二年（319年）] \eventanchor{JH-LEDGER-033-2}{前赵·“刘曜在长安改汉国号为赵，以关中为中心重建政权身份”的命令、联盟或地方响应}
  \item[光初二年（319年）] \eventanchor{JH-LEDGER-033-3}{前赵·刘曜在长安改汉国号为赵，以关中为中心重建政权身份}
  \item[光初二年（319年）] \eventanchor{JH-LEDGER-033-4}{前赵·“刘曜在长安改汉国号为赵，以关中为中心重建政权身份”的执行中的空间与人员处置}
  \item[光初二年（319年）] \eventanchor{JH-LEDGER-033-5}{前赵·“刘曜在长安改汉国号为赵，以关中为中心重建政权身份”的后续制度或军政调整}
  \item[光初二年（319年）] \eventanchor{JH-LEDGER-033-6}{前赵·“刘曜在长安改汉国号为赵，以关中为中心重建政权身份”的异源材料与影响范围的复核}
  \item[大兴三年（320年）] \eventanchor{JH-LEDGER-034-1}{东晋与荆湘地方集团·“杜弢领导的荆湘叛乱被王敦等镇压，东晋重新取得长江中游部分军政控制”的前期动员与资源准备}
  \item[大兴三年（320年）] \eventanchor{JH-LEDGER-034-2}{东晋与荆湘地方集团·“杜弢领导的荆湘叛乱被王敦等镇压，东晋重新取得长江中游部分军政控制”的命令、联盟或地方响应}
  \item[大兴三年（320年）] \eventanchor{JH-LEDGER-034-3}{东晋与荆湘地方集团·杜弢领导的荆湘叛乱被王敦等镇压，东晋重新取得长江中游部分军政控制}
  \item[大兴三年（320年）] \eventanchor{JH-LEDGER-034-4}{东晋与荆湘地方集团·“杜弢领导的荆湘叛乱被王敦等镇压，东晋重新取得长江中游部分军政控制”的执行中的空间与人员处置}
  \item[大兴三年（320年）] \eventanchor{JH-LEDGER-034-5}{东晋与荆湘地方集团·“杜弢领导的荆湘叛乱被王敦等镇压，东晋重新取得长江中游部分军政控制”的后续制度或军政调整}
  \item[大兴三年（320年）] \eventanchor{JH-LEDGER-034-6}{东晋与荆湘地方集团·“杜弢领导的荆湘叛乱被王敦等镇压，东晋重新取得长江中游部分军政控制”的异源材料与影响范围的复核}
  \item[446年] \eventanchor{NS-LEDGER-001-1}{北魏·“太武帝禁佛”的前期动员与资源准备}
  \item[446年] \eventanchor{NS-LEDGER-001-2}{北魏·“太武帝禁佛”的命令、联盟或地方响应}
  \item[446年] \eventanchor{NS-LEDGER-001-3}{北魏·太武帝禁佛}
  \item[446年] \eventanchor{NS-LEDGER-001-4}{北魏·“太武帝禁佛”的执行中的空间与人员处置}
  \item[446年] \eventanchor{NS-LEDGER-001-5}{北魏·“太武帝禁佛”的后续制度或军政调整}
  \item[446年] \eventanchor{NS-LEDGER-001-6}{北魏·“太武帝禁佛”的异源材料与影响范围的复核}
  \item[452年] \eventanchor{NS-LEDGER-002-1}{北魏·“文成帝即位后恢复佛教”的前期动员与资源准备}
  \item[452年] \eventanchor{NS-LEDGER-002-2}{北魏·“文成帝即位后恢复佛教”的命令、联盟或地方响应}
  \item[452年] \eventanchor{NS-LEDGER-002-3}{北魏·文成帝即位后恢复佛教}
  \item[452年] \eventanchor{NS-LEDGER-002-4}{北魏·“文成帝即位后恢复佛教”的执行中的空间与人员处置}
  \item[452年] \eventanchor{NS-LEDGER-002-5}{北魏·“文成帝即位后恢复佛教”的后续制度或军政调整}
  \item[452年] \eventanchor{NS-LEDGER-002-6}{北魏·“文成帝即位后恢复佛教”的异源材料与影响范围的复核}
  \item[453年] \eventanchor{NS-LEDGER-003-1}{刘宋·“太子劭弑文帝”的前期动员与资源准备}
  \item[453年] \eventanchor{NS-LEDGER-003-2}{刘宋·“太子劭弑文帝”的命令、联盟或地方响应}
  \item[453年] \eventanchor{NS-LEDGER-003-3}{刘宋·太子劭弑文帝}
  \item[453年] \eventanchor{NS-LEDGER-003-4}{刘宋·“太子劭弑文帝”的执行中的空间与人员处置}
  \item[453年] \eventanchor{NS-LEDGER-003-5}{刘宋·“太子劭弑文帝”的后续制度或军政调整}
  \item[453年] \eventanchor{NS-LEDGER-003-6}{刘宋·“太子劭弑文帝”的异源材料与影响范围的复核}
  \item[454年] \eventanchor{NS-LEDGER-004-1}{刘宋·“孝武帝平定刘劭”的前期动员与资源准备}
  \item[454年] \eventanchor{NS-LEDGER-004-2}{刘宋·“孝武帝平定刘劭”的命令、联盟或地方响应}
  \item[454年] \eventanchor{NS-LEDGER-004-3}{刘宋·孝武帝平定刘劭}
  \item[454年] \eventanchor{NS-LEDGER-004-4}{刘宋·“孝武帝平定刘劭”的执行中的空间与人员处置}
  \item[454年] \eventanchor{NS-LEDGER-004-5}{刘宋·“孝武帝平定刘劭”的后续制度或军政调整}
  \item[454年] \eventanchor{NS-LEDGER-004-6}{刘宋·“孝武帝平定刘劭”的异源材料与影响范围的复核}
  \item[466年] \eventanchor{NS-LEDGER-005-1}{刘宋·“子勋之乱”的前期动员与资源准备}
  \item[466年] \eventanchor{NS-LEDGER-005-2}{刘宋·“子勋之乱”的命令、联盟或地方响应}
  \item[466年] \eventanchor{NS-LEDGER-005-3}{刘宋·子勋之乱}
  \item[466年] \eventanchor{NS-LEDGER-005-4}{刘宋·“子勋之乱”的执行中的空间与人员处置}
  \item[466年] \eventanchor{NS-LEDGER-005-5}{刘宋·“子勋之乱”的后续制度或军政调整}
  \item[466年] \eventanchor{NS-LEDGER-005-6}{刘宋·“子勋之乱”的异源材料与影响范围的复核}
  \item[477年] \eventanchor{NS-LEDGER-006-1}{刘宋·“萧道成拥立顺帝”的前期动员与资源准备}
  \item[477年] \eventanchor{NS-LEDGER-006-2}{刘宋·“萧道成拥立顺帝”的命令、联盟或地方响应}
  \item[477年] \eventanchor{NS-LEDGER-006-3}{刘宋·萧道成拥立顺帝}
  \item[477年] \eventanchor{NS-LEDGER-006-4}{刘宋·“萧道成拥立顺帝”的执行中的空间与人员处置}
  \item[477年] \eventanchor{NS-LEDGER-006-5}{刘宋·“萧道成拥立顺帝”的后续制度或军政调整}
  \item[477年] \eventanchor{NS-LEDGER-006-6}{刘宋·“萧道成拥立顺帝”的异源材料与影响范围的复核}
  \item[479年] \eventanchor{NS-LEDGER-007-1}{南齐·“萧道成受禅建齐”的前期动员与资源准备}
  \item[479年] \eventanchor{NS-LEDGER-007-2}{南齐·“萧道成受禅建齐”的命令、联盟或地方响应}
  \item[479年] \eventanchor{NS-LEDGER-007-3}{南齐·萧道成受禅建齐}
  \item[479年] \eventanchor{NS-LEDGER-007-4}{南齐·“萧道成受禅建齐”的执行中的空间与人员处置}
  \item[479年] \eventanchor{NS-LEDGER-007-5}{南齐·“萧道成受禅建齐”的后续制度或军政调整}
  \item[479年] \eventanchor{NS-LEDGER-007-6}{南齐·“萧道成受禅建齐”的异源材料与影响范围的复核}
  \item[483年] \eventanchor{NS-LEDGER-008-1}{南齐·“齐武帝即位”的前期动员与资源准备}
  \item[483年] \eventanchor{NS-LEDGER-008-2}{南齐·“齐武帝即位”的命令、联盟或地方响应}
  \item[483年] \eventanchor{NS-LEDGER-008-3}{南齐·齐武帝即位}
  \item[483年] \eventanchor{NS-LEDGER-008-4}{南齐·“齐武帝即位”的执行中的空间与人员处置}
  \item[483年] \eventanchor{NS-LEDGER-008-5}{南齐·“齐武帝即位”的后续制度或军政调整}
  \item[483年] \eventanchor{NS-LEDGER-008-6}{南齐·“齐武帝即位”的异源材料与影响范围的复核}
  \item[494年] \eventanchor{NS-LEDGER-009-1}{南齐·“齐明帝即位”的前期动员与资源准备}
  \item[494年] \eventanchor{NS-LEDGER-009-2}{南齐·“齐明帝即位”的命令、联盟或地方响应}
  \item[494年] \eventanchor{NS-LEDGER-009-3}{南齐·齐明帝即位}
  \item[494年] \eventanchor{NS-LEDGER-009-4}{南齐·“齐明帝即位”的执行中的空间与人员处置}
  \item[494年] \eventanchor{NS-LEDGER-009-5}{南齐·“齐明帝即位”的后续制度或军政调整}
  \item[494年] \eventanchor{NS-LEDGER-009-6}{南齐·“齐明帝即位”的异源材料与影响范围的复核}
  \item[501年] \eventanchor{NS-LEDGER-010-1}{南齐与梁·“萧衍起兵”的前期动员与资源准备}
  \item[501年] \eventanchor{NS-LEDGER-010-2}{南齐与梁·“萧衍起兵”的命令、联盟或地方响应}
  \item[501年] \eventanchor{NS-LEDGER-010-3}{南齐与梁·萧衍起兵}
  \item[501年] \eventanchor{NS-LEDGER-010-4}{南齐与梁·“萧衍起兵”的执行中的空间与人员处置}
  \item[501年] \eventanchor{NS-LEDGER-010-5}{南齐与梁·“萧衍起兵”的后续制度或军政调整}
  \item[501年] \eventanchor{NS-LEDGER-010-6}{南齐与梁·“萧衍起兵”的异源材料与影响范围的复核}
  \item[502年] \eventanchor{NS-LEDGER-011-1}{梁·“萧衍建梁”的前期动员与资源准备}
  \item[502年] \eventanchor{NS-LEDGER-011-2}{梁·“萧衍建梁”的命令、联盟或地方响应}
  \item[502年] \eventanchor{NS-LEDGER-011-3}{梁·萧衍建梁}
  \item[502年] \eventanchor{NS-LEDGER-011-4}{梁·“萧衍建梁”的执行中的空间与人员处置}
  \item[502年] \eventanchor{NS-LEDGER-011-5}{梁·“萧衍建梁”的后续制度或军政调整}
  \item[502年] \eventanchor{NS-LEDGER-011-6}{梁·“萧衍建梁”的异源材料与影响范围的复核}
  \item[504年] \eventanchor{NS-LEDGER-012-1}{梁·“梁武帝公开崇佛政策”的前期动员与资源准备}
  \item[504年] \eventanchor{NS-LEDGER-012-2}{梁·“梁武帝公开崇佛政策”的命令、联盟或地方响应}
  \item[504年] \eventanchor{NS-LEDGER-012-3}{梁·梁武帝公开崇佛政策}
  \item[504年] \eventanchor{NS-LEDGER-012-4}{梁·“梁武帝公开崇佛政策”的执行中的空间与人员处置}
  \item[504年] \eventanchor{NS-LEDGER-012-5}{梁·“梁武帝公开崇佛政策”的后续制度或军政调整}
  \item[504年] \eventanchor{NS-LEDGER-012-6}{梁·“梁武帝公开崇佛政策”的异源材料与影响范围的复核}
  \item[523年] \eventanchor{NS-LEDGER-013-1}{北魏·“六镇起义爆发”的前期动员与资源准备}
  \item[523年] \eventanchor{NS-LEDGER-013-2}{北魏·“六镇起义爆发”的命令、联盟或地方响应}
  \item[523年] \eventanchor{NS-LEDGER-013-3}{北魏·六镇起义爆发}
  \item[523年] \eventanchor{NS-LEDGER-013-4}{北魏·“六镇起义爆发”的执行中的空间与人员处置}
  \item[523年] \eventanchor{NS-LEDGER-013-5}{北魏·“六镇起义爆发”的后续制度或军政调整}
  \item[523年] \eventanchor{NS-LEDGER-013-6}{北魏·“六镇起义爆发”的异源材料与影响范围的复核}
  \item[528年] \eventanchor{NS-LEDGER-014-1}{北魏·“孝明帝死，尔朱荣入洛”的前期动员与资源准备}
  \item[528年] \eventanchor{NS-LEDGER-014-2}{北魏·“孝明帝死，尔朱荣入洛”的命令、联盟或地方响应}
  \item[528年] \eventanchor{NS-LEDGER-014-3}{北魏·孝明帝死，尔朱荣入洛}
  \item[528年] \eventanchor{NS-LEDGER-014-4}{北魏·“孝明帝死，尔朱荣入洛”的执行中的空间与人员处置}
  \item[528年] \eventanchor{NS-LEDGER-014-5}{北魏·“孝明帝死，尔朱荣入洛”的后续制度或军政调整}
  \item[528年] \eventanchor{NS-LEDGER-014-6}{北魏·“孝明帝死，尔朱荣入洛”的异源材料与影响范围的复核}
  \item[530年] \eventanchor{NS-LEDGER-015-1}{北魏·“孝庄帝杀尔朱荣”的前期动员与资源准备}
  \item[530年] \eventanchor{NS-LEDGER-015-2}{北魏·“孝庄帝杀尔朱荣”的命令、联盟或地方响应}
  \item[530年] \eventanchor{NS-LEDGER-015-3}{北魏·孝庄帝杀尔朱荣}
  \item[530年] \eventanchor{NS-LEDGER-015-4}{北魏·“孝庄帝杀尔朱荣”的执行中的空间与人员处置}
  \item[530年] \eventanchor{NS-LEDGER-015-5}{北魏·“孝庄帝杀尔朱荣”的后续制度或军政调整}
  \item[530年] \eventanchor{NS-LEDGER-015-6}{北魏·“孝庄帝杀尔朱荣”的异源材料与影响范围的复核}
  \item[532年] \eventanchor{NS-LEDGER-016-1}{北魏·“高欢控制东部朝廷”的前期动员与资源准备}
  \item[532年] \eventanchor{NS-LEDGER-016-2}{北魏·“高欢控制东部朝廷”的命令、联盟或地方响应}
  \item[532年] \eventanchor{NS-LEDGER-016-3}{北魏·高欢控制东部朝廷}
  \item[532年] \eventanchor{NS-LEDGER-016-4}{北魏·“高欢控制东部朝廷”的执行中的空间与人员处置}
  \item[532年] \eventanchor{NS-LEDGER-016-5}{北魏·“高欢控制东部朝廷”的后续制度或军政调整}
  \item[532年] \eventanchor{NS-LEDGER-016-6}{北魏·“高欢控制东部朝廷”的异源材料与影响范围的复核}
  \item[534年] \eventanchor{NS-LEDGER-017-1}{东魏与西魏·“北魏分裂为东魏西魏”的前期动员与资源准备}
  \item[534年] \eventanchor{NS-LEDGER-017-2}{东魏与西魏·“北魏分裂为东魏西魏”的命令、联盟或地方响应}
  \item[534年] \eventanchor{NS-LEDGER-017-3}{东魏与西魏·北魏分裂为东魏西魏}
  \item[534年] \eventanchor{NS-LEDGER-017-4}{东魏与西魏·“北魏分裂为东魏西魏”的执行中的空间与人员处置}
  \item[534年] \eventanchor{NS-LEDGER-017-5}{东魏与西魏·“北魏分裂为东魏西魏”的后续制度或军政调整}
  \item[534年] \eventanchor{NS-LEDGER-017-6}{东魏与西魏·“北魏分裂为东魏西魏”的异源材料与影响范围的复核}
  \item[535年] \eventanchor{NS-LEDGER-018-1}{东魏·“高欢经营邺城军政”的前期动员与资源准备}
  \item[535年] \eventanchor{NS-LEDGER-018-2}{东魏·“高欢经营邺城军政”的命令、联盟或地方响应}
  \item[535年] \eventanchor{NS-LEDGER-018-3}{东魏·高欢经营邺城军政}
  \item[535年] \eventanchor{NS-LEDGER-018-4}{东魏·“高欢经营邺城军政”的执行中的空间与人员处置}
  \item[535年] \eventanchor{NS-LEDGER-018-5}{东魏·“高欢经营邺城军政”的后续制度或军政调整}
  \item[535年] \eventanchor{NS-LEDGER-018-6}{东魏·“高欢经营邺城军政”的异源材料与影响范围的复核}
  \item[535年] \eventanchor{NS-LEDGER-019-1}{西魏·“宇文泰控制长安”的前期动员与资源准备}
  \item[535年] \eventanchor{NS-LEDGER-019-2}{西魏·“宇文泰控制长安”的命令、联盟或地方响应}
  \item[535年] \eventanchor{NS-LEDGER-019-3}{西魏·宇文泰控制长安}
  \item[535年] \eventanchor{NS-LEDGER-019-4}{西魏·“宇文泰控制长安”的执行中的空间与人员处置}
  \item[535年] \eventanchor{NS-LEDGER-019-5}{西魏·“宇文泰控制长安”的后续制度或军政调整}
  \item[535年] \eventanchor{NS-LEDGER-019-6}{西魏·“宇文泰控制长安”的异源材料与影响范围的复核}
  \item[547年] \eventanchor{NS-LEDGER-020-1}{东魏与梁·“侯景叛高欢南投梁”的前期动员与资源准备}
  \item[547年] \eventanchor{NS-LEDGER-020-2}{东魏与梁·“侯景叛高欢南投梁”的命令、联盟或地方响应}
  \item[547年] \eventanchor{NS-LEDGER-020-3}{东魏与梁·侯景叛高欢南投梁}
  \item[547年] \eventanchor{NS-LEDGER-020-4}{东魏与梁·“侯景叛高欢南投梁”的执行中的空间与人员处置}
  \item[547年] \eventanchor{NS-LEDGER-020-5}{东魏与梁·“侯景叛高欢南投梁”的后续制度或军政调整}
  \item[547年] \eventanchor{NS-LEDGER-020-6}{东魏与梁·“侯景叛高欢南投梁”的异源材料与影响范围的复核}
  \item[548年] \eventanchor{NS-LEDGER-021-1}{梁与侯景·“侯景渡江反梁”的前期动员与资源准备}
  \item[548年] \eventanchor{NS-LEDGER-021-2}{梁与侯景·“侯景渡江反梁”的命令、联盟或地方响应}
  \item[548年] \eventanchor{NS-LEDGER-021-3}{梁与侯景·侯景渡江反梁}
  \item[548年] \eventanchor{NS-LEDGER-021-4}{梁与侯景·“侯景渡江反梁”的执行中的空间与人员处置}
  \item[548年] \eventanchor{NS-LEDGER-021-5}{梁与侯景·“侯景渡江反梁”的后续制度或军政调整}
  \item[548年] \eventanchor{NS-LEDGER-021-6}{梁与侯景·“侯景渡江反梁”的异源材料与影响范围的复核}
  \item[549年] \eventanchor{NS-LEDGER-022-1}{梁·“侯景攻陷建康”的前期动员与资源准备}
  \item[549年] \eventanchor{NS-LEDGER-022-2}{梁·“侯景攻陷建康”的命令、联盟或地方响应}
  \item[549年] \eventanchor{NS-LEDGER-022-3}{梁·侯景攻陷建康}
  \item[549年] \eventanchor{NS-LEDGER-022-4}{梁·“侯景攻陷建康”的执行中的空间与人员处置}
  \item[549年] \eventanchor{NS-LEDGER-022-5}{梁·“侯景攻陷建康”的后续制度或军政调整}
  \item[549年] \eventanchor{NS-LEDGER-022-6}{梁·“侯景攻陷建康”的异源材料与影响范围的复核}
  \item[552年] \eventanchor{NS-LEDGER-023-1}{梁·“王僧辩陈霸先平侯景”的前期动员与资源准备}
  \item[552年] \eventanchor{NS-LEDGER-023-2}{梁·“王僧辩陈霸先平侯景”的命令、联盟或地方响应}
  \item[552年] \eventanchor{NS-LEDGER-023-3}{梁·王僧辩陈霸先平侯景}
  \item[552年] \eventanchor{NS-LEDGER-023-4}{梁·“王僧辩陈霸先平侯景”的执行中的空间与人员处置}
  \item[552年] \eventanchor{NS-LEDGER-023-5}{梁·“王僧辩陈霸先平侯景”的后续制度或军政调整}
  \item[552年] \eventanchor{NS-LEDGER-023-6}{梁·“王僧辩陈霸先平侯景”的异源材料与影响范围的复核}
  \item[554年] \eventanchor{NS-LEDGER-024-1}{西魏与梁·“西魏攻取江陵”的前期动员与资源准备}
  \item[554年] \eventanchor{NS-LEDGER-024-2}{西魏与梁·“西魏攻取江陵”的命令、联盟或地方响应}
  \item[554年] \eventanchor{NS-LEDGER-024-3}{西魏与梁·西魏攻取江陵}
  \item[554年] \eventanchor{NS-LEDGER-024-4}{西魏与梁·“西魏攻取江陵”的执行中的空间与人员处置}
  \item[554年] \eventanchor{NS-LEDGER-024-5}{西魏与梁·“西魏攻取江陵”的后续制度或军政调整}
  \item[554年] \eventanchor{NS-LEDGER-024-6}{西魏与梁·“西魏攻取江陵”的异源材料与影响范围的复核}
  \item[557年] \eventanchor{NS-LEDGER-025-1}{陈·“陈霸先建陈”的前期动员与资源准备}
  \item[557年] \eventanchor{NS-LEDGER-025-2}{陈·“陈霸先建陈”的命令、联盟或地方响应}
  \item[557年] \eventanchor{NS-LEDGER-025-3}{陈·陈霸先建陈}
  \item[557年] \eventanchor{NS-LEDGER-025-4}{陈·“陈霸先建陈”的执行中的空间与人员处置}
  \item[557年] \eventanchor{NS-LEDGER-025-5}{陈·“陈霸先建陈”的后续制度或军政调整}
  \item[557年] \eventanchor{NS-LEDGER-025-6}{陈·“陈霸先建陈”的异源材料与影响范围的复核}
  \item[557年] \eventanchor{NS-LEDGER-026-1}{北周·“宇文觉建北周”的前期动员与资源准备}
  \item[557年] \eventanchor{NS-LEDGER-026-2}{北周·“宇文觉建北周”的命令、联盟或地方响应}
  \item[557年] \eventanchor{NS-LEDGER-026-3}{北周·宇文觉建北周}
  \item[557年] \eventanchor{NS-LEDGER-026-4}{北周·“宇文觉建北周”的执行中的空间与人员处置}
  \item[557年] \eventanchor{NS-LEDGER-026-5}{北周·“宇文觉建北周”的后续制度或军政调整}
  \item[557年] \eventanchor{NS-LEDGER-026-6}{北周·“宇文觉建北周”的异源材料与影响范围的复核}
  \item[565年] \eventanchor{NS-LEDGER-027-1}{北齐·“后主继位”的前期动员与资源准备}
  \item[565年] \eventanchor{NS-LEDGER-027-2}{北齐·“后主继位”的命令、联盟或地方响应}
  \item[565年] \eventanchor{NS-LEDGER-027-3}{北齐·后主继位}
  \item[565年] \eventanchor{NS-LEDGER-027-4}{北齐·“后主继位”的执行中的空间与人员处置}
  \item[565年] \eventanchor{NS-LEDGER-027-5}{北齐·“后主继位”的后续制度或军政调整}
  \item[565年] \eventanchor{NS-LEDGER-027-6}{北齐·“后主继位”的异源材料与影响范围的复核}
  \item[572年] \eventanchor{NS-LEDGER-028-1}{北周·“周武帝诛宇文护”的前期动员与资源准备}
  \item[572年] \eventanchor{NS-LEDGER-028-2}{北周·“周武帝诛宇文护”的命令、联盟或地方响应}
  \item[572年] \eventanchor{NS-LEDGER-028-3}{北周·周武帝诛宇文护}
  \item[572年] \eventanchor{NS-LEDGER-028-4}{北周·“周武帝诛宇文护”的执行中的空间与人员处置}
  \item[572年] \eventanchor{NS-LEDGER-028-5}{北周·“周武帝诛宇文护”的后续制度或军政调整}
  \item[572年] \eventanchor{NS-LEDGER-028-6}{北周·“周武帝诛宇文护”的异源材料与影响范围的复核}
  \item[574年] \eventanchor{NS-LEDGER-029-1}{北周·“周武帝禁佛道”的前期动员与资源准备}
  \item[574年] \eventanchor{NS-LEDGER-029-2}{北周·“周武帝禁佛道”的命令、联盟或地方响应}
  \item[574年] \eventanchor{NS-LEDGER-029-3}{北周·周武帝禁佛道}
  \item[574年] \eventanchor{NS-LEDGER-029-4}{北周·“周武帝禁佛道”的执行中的空间与人员处置}
  \item[574年] \eventanchor{NS-LEDGER-029-5}{北周·“周武帝禁佛道”的后续制度或军政调整}
  \item[574年] \eventanchor{NS-LEDGER-029-6}{北周·“周武帝禁佛道”的异源材料与影响范围的复核}
  \item[577年] \eventanchor{NS-LEDGER-030-1}{北周与北齐·“北周灭北齐”的前期动员与资源准备}
  \item[577年] \eventanchor{NS-LEDGER-030-2}{北周与北齐·“北周灭北齐”的命令、联盟或地方响应}
  \item[577年] \eventanchor{NS-LEDGER-030-3}{北周与北齐·北周灭北齐}
  \item[577年] \eventanchor{NS-LEDGER-030-4}{北周与北齐·“北周灭北齐”的执行中的空间与人员处置}
  \item[577年] \eventanchor{NS-LEDGER-030-5}{北周与北齐·“北周灭北齐”的后续制度或军政调整}
  \item[577年] \eventanchor{NS-LEDGER-030-6}{北周与北齐·“北周灭北齐”的异源材料与影响范围的复核}
  \item[578年] \eventanchor{NS-LEDGER-031-1}{北周·“周武帝去世”的前期动员与资源准备}
  \item[578年] \eventanchor{NS-LEDGER-031-2}{北周·“周武帝去世”的命令、联盟或地方响应}
  \item[578年] \eventanchor{NS-LEDGER-031-3}{北周·周武帝去世}
  \item[578年] \eventanchor{NS-LEDGER-031-4}{北周·“周武帝去世”的执行中的空间与人员处置}
  \item[578年] \eventanchor{NS-LEDGER-031-5}{北周·“周武帝去世”的后续制度或军政调整}
  \item[578年] \eventanchor{NS-LEDGER-031-6}{北周·“周武帝去世”的异源材料与影响范围的复核}
  \item[580年] \eventanchor{NS-LEDGER-032-1}{北周·“杨坚辅政”的前期动员与资源准备}
  \item[580年] \eventanchor{NS-LEDGER-032-2}{北周·“杨坚辅政”的命令、联盟或地方响应}
  \item[580年] \eventanchor{NS-LEDGER-032-3}{北周·杨坚辅政}
  \item[580年] \eventanchor{NS-LEDGER-032-4}{北周·“杨坚辅政”的执行中的空间与人员处置}
  \item[580年] \eventanchor{NS-LEDGER-032-5}{北周·“杨坚辅政”的后续制度或军政调整}
  \item[580年] \eventanchor{NS-LEDGER-032-6}{北周·“杨坚辅政”的异源材料与影响范围的复核}
  \item[581年] \eventanchor{NS-LEDGER-033-1}{隋·“杨坚受禅建隋”的前期动员与资源准备}
  \item[581年] \eventanchor{NS-LEDGER-033-2}{隋·“杨坚受禅建隋”的命令、联盟或地方响应}
  \item[581年] \eventanchor{NS-LEDGER-033-3}{隋·杨坚受禅建隋}
  \item[581年] \eventanchor{NS-LEDGER-033-4}{隋·“杨坚受禅建隋”的执行中的空间与人员处置}
  \item[581年] \eventanchor{NS-LEDGER-033-5}{隋·“杨坚受禅建隋”的后续制度或军政调整}
  \item[581年] \eventanchor{NS-LEDGER-033-6}{隋·“杨坚受禅建隋”的异源材料与影响范围的复核}
  \item[583年] \eventanchor{NS-LEDGER-034-1}{隋·“隋整顿北方军政与户籍”的前期动员与资源准备}
  \item[583年] \eventanchor{NS-LEDGER-034-2}{隋·“隋整顿北方军政与户籍”的命令、联盟或地方响应}
  \item[583年] \eventanchor{NS-LEDGER-034-3}{隋·隋整顿北方军政与户籍}
  \item[583年] \eventanchor{NS-LEDGER-034-4}{隋·“隋整顿北方军政与户籍”的执行中的空间与人员处置}
  \item[583年] \eventanchor{NS-LEDGER-034-5}{隋·“隋整顿北方军政与户籍”的后续制度或军政调整}
  \item[583年] \eventanchor{NS-LEDGER-034-6}{隋·“隋整顿北方军政与户籍”的异源材料与影响范围的复核}
  \item[588年] \eventanchor{NS-LEDGER-035-1}{隋与陈·“隋军南下伐陈”的前期动员与资源准备}
  \item[588年] \eventanchor{NS-LEDGER-035-2}{隋与陈·“隋军南下伐陈”的命令、联盟或地方响应}
  \item[588年] \eventanchor{NS-LEDGER-035-3}{隋与陈·隋军南下伐陈}
  \item[588年] \eventanchor{NS-LEDGER-035-4}{隋与陈·“隋军南下伐陈”的执行中的空间与人员处置}
  \item[588年] \eventanchor{NS-LEDGER-035-5}{隋与陈·“隋军南下伐陈”的后续制度或军政调整}
  \item[588年] \eventanchor{NS-LEDGER-035-6}{隋与陈·“隋军南下伐陈”的异源材料与影响范围的复核}
  \item[589年] \eventanchor{NS-LEDGER-036-1}{隋与陈·“建康失陷，陈亡”的前期动员与资源准备}
  \item[589年] \eventanchor{NS-LEDGER-036-2}{隋与陈·“建康失陷，陈亡”的命令、联盟或地方响应}
  \item[589年] \eventanchor{NS-LEDGER-036-3}{隋与陈·建康失陷，陈亡}
  \item[589年] \eventanchor{NS-LEDGER-036-4}{隋与陈·“建康失陷，陈亡”的执行中的空间与人员处置}
  \item[589年] \eventanchor{NS-LEDGER-036-5}{隋与陈·“建康失陷，陈亡”的后续制度或军政调整}
  \item[589年] \eventanchor{NS-LEDGER-036-6}{隋与陈·“建康失陷，陈亡”的异源材料与影响范围的复核}
  \item[450年] \eventanchor{NS-LEDGER-037-1}{北魏与刘宋·“北魏南侵至江淮”的前期动员与资源准备}
  \item[450年] \eventanchor{NS-LEDGER-037-2}{北魏与刘宋·“北魏南侵至江淮”的命令、联盟或地方响应}
  \item[450年] \eventanchor{NS-LEDGER-037-3}{北魏与刘宋·北魏南侵至江淮}
  \item[450年] \eventanchor{NS-LEDGER-037-4}{北魏与刘宋·“北魏南侵至江淮”的执行中的空间与人员处置}
  \item[450年] \eventanchor{NS-LEDGER-037-5}{北魏与刘宋·“北魏南侵至江淮”的后续制度或军政调整}
  \item[450年] \eventanchor{NS-LEDGER-037-6}{北魏与刘宋·“北魏南侵至江淮”的异源材料与影响范围的复核}
  \item[471年] \eventanchor{NS-LEDGER-038-1}{北魏·“献文帝禅位，孝文帝继位”的前期动员与资源准备}
  \item[471年] \eventanchor{NS-LEDGER-038-2}{北魏·“献文帝禅位，孝文帝继位”的命令、联盟或地方响应}
  \item[471年] \eventanchor{NS-LEDGER-038-3}{北魏·献文帝禅位，孝文帝继位}
  \item[471年] \eventanchor{NS-LEDGER-038-4}{北魏·“献文帝禅位，孝文帝继位”的执行中的空间与人员处置}
  \item[471年] \eventanchor{NS-LEDGER-038-5}{北魏·“献文帝禅位，孝文帝继位”的后续制度或军政调整}
  \item[471年] \eventanchor{NS-LEDGER-038-6}{北魏·“献文帝禅位，孝文帝继位”的异源材料与影响范围的复核}
  \item[494年] \eventanchor{NS-LEDGER-039-1}{北魏·“北魏迁都洛阳”的前期动员与资源准备}
  \item[494年] \eventanchor{NS-LEDGER-039-2}{北魏·“北魏迁都洛阳”的命令、联盟或地方响应}
  \item[494年] \eventanchor{NS-LEDGER-039-3}{北魏·北魏迁都洛阳}
  \item[494年] \eventanchor{NS-LEDGER-039-4}{北魏·“北魏迁都洛阳”的执行中的空间与人员处置}
  \item[494年] \eventanchor{NS-LEDGER-039-5}{北魏·“北魏迁都洛阳”的后续制度或军政调整}
  \item[494年] \eventanchor{NS-LEDGER-039-6}{北魏·“北魏迁都洛阳”的异源材料与影响范围的复核}
  \item[495年] \eventanchor{NS-LEDGER-040-1}{北魏·“孝文帝改革姓氏服制”的前期动员与资源准备}
  \item[495年] \eventanchor{NS-LEDGER-040-2}{北魏·“孝文帝改革姓氏服制”的命令、联盟或地方响应}
  \item[495年] \eventanchor{NS-LEDGER-040-3}{北魏·孝文帝改革姓氏服制}
  \item[495年] \eventanchor{NS-LEDGER-040-4}{北魏·“孝文帝改革姓氏服制”的执行中的空间与人员处置}
  \item[495年] \eventanchor{NS-LEDGER-040-5}{北魏·“孝文帝改革姓氏服制”的后续制度或军政调整}
  \item[495年] \eventanchor{NS-LEDGER-040-6}{北魏·“孝文帝改革姓氏服制”的异源材料与影响范围的复核}
  \item[440年] \eventanchor{NSL-440-589-001}{北魏·北魏在河西旧北凉地区配置镇戍并接收降民}
  \item[约440年] \eventanchor{NSL-440-589-002}{北凉遗众与高昌·沮渠无讳率北凉遗众西趋高昌，延续河西政权的人口与政治网络}
  \item[441年] \eventanchor{NSL-440-589-003}{北魏与柔然·北魏北边诸镇继续防备柔然南下与掠夺}
  \item[442年] \eventanchor{NSL-440-589-004}{高昌·沮渠无讳在高昌建立统治并吸纳河西随从}
  \item[443年] \eventanchor{NSL-440-589-005}{北魏与柔然·太武帝亲率北征，分道攻击柔然部众}
  \item[443年] \eventanchor{NSL-440-589-006}{北魏·北魏将部分高车、柔然降附人群安置于北部边地}
  \item[444年] \eventanchor{NSL-440-589-007}{北魏·太武朝继续整饬北边军镇与畜牧供给}
  \item[445年] \eventanchor{NSL-440-589-008}{北魏·盖吴在杏城一带起兵，联络关中民众反魏}
  \item[445年] \eventanchor{NSL-440-589-009}{北魏·北魏以崔浩等参与军政筹划，调兵镇压盖吴集团}
  \item[445年] \eventanchor{NSL-440-589-010}{北魏·长安僧人玄高等被卷入盖吴事件后的审讯与处置}
  \item[446年] \eventanchor{NSL-440-589-011}{北魏·太武帝平定盖吴余部，关中起事网络被武力拆解}
  \item[446年] \eventanchor{NSL-440-589-012}{北魏·太武帝下令毁寺、禁佛，僧尼被迫还俗或逃散}
  \item[446年] \eventanchor{NSL-440-589-013}{北魏·北魏禁佛令的地方执行程度因军镇与州郡条件而不一}
  \item[447年] \eventanchor{NSL-440-589-014}{北魏·太武朝继续处理关中降附者、流民与被毁寺院遗产}
  \item[448年] \eventanchor{NSL-440-589-015}{北魏·新天师道领袖寇谦之去世，北魏宫廷的道教政治联盟失去关键人物}
  \item[448年] \eventanchor{NSL-440-589-016}{北魏与柔然·北魏北方军镇持续报告柔然活动并维持边防巡察}
  \item[449年] \eventanchor{NSL-440-589-017}{北魏·太武朝在关中、河西继续以镇将和州郡官分担边地治理}
  \item[450年] \eventanchor{NSL-440-589-018}{刘宋·文帝发动元嘉北伐，宋军自淮、泗方向北进}
  \item[450年] \eventanchor{NSL-440-589-019}{北魏与刘宋·太武帝反击南朝，魏军越淮北进并冲击彭城、瓜步方向}
  \item[450年] \eventanchor{NSL-440-589-020}{刘宋·宋廷在北伐受挫后动员沿江防御与都城戒严}
  \item[451年] \eventanchor{NSL-440-589-021}{北魏与刘宋·魏军自江北撤回，宋魏战线重新收缩至淮、泗一带}
  \item[451年] \eventanchor{NSL-440-589-022}{刘宋·元嘉北伐后宋廷检讨将领失利与边防部署}
  \item[452年] \eventanchor{NSL-440-589-023}{北魏·太武帝被宗爱杀害，南安王余短暂即位后亦被杀}
  \item[452年] \eventanchor{NSL-440-589-024}{北魏·文成帝即位，冯太后与大臣集团协助重建朝廷秩序}
  \item[452年] \eventanchor{NSL-440-589-025}{北魏·文成帝恢复佛教，允许僧尼与寺院重新活动}
  \item[453年] \eventanchor{NSL-440-589-026}{刘宋·太子刘劭弑杀文帝，自立为帝}
  \item[453年] \eventanchor{NSL-440-589-027}{刘宋·武陵王刘骏在江州起兵，联结地方军府反刘劭}
  \item[453年] \eventanchor{NSL-440-589-028}{刘宋·刘骏攻入建康，处死刘劭并即帝位，是为孝武帝}
  \item[454年] \eventanchor{NSL-440-589-029}{刘宋·孝武帝平定南郡王刘义宣、臧质等的反叛}
  \item[454年] \eventanchor{NSL-440-589-030}{刘宋·刘宋朝廷在内战后加强对宗室王镇与州郡兵权的约束}
  \item[455年] \eventanchor{NSL-440-589-031}{北魏·文成朝继续修复太武末年受损的宫廷与行政秩序}
  \item[456年] \eventanchor{NSL-440-589-032}{北魏·北魏在平城周边整饬宫城、苑囿与供给体系}
  \item[457年] \eventanchor{NSL-440-589-033}{刘宋·孝武帝以皇族封授和州镇调整安置内战后功臣与宗室}
  \item[458年] \eventanchor{NSL-440-589-034}{北魏与柔然·柔然与北魏在漠南牧地和边镇周边反复争夺}
  \item[459年] \eventanchor{NSL-440-589-035}{刘宋·宋廷继续经营淮北沿线城戍并整修江防}
  \item[460年] \eventanchor{NSL-440-589-036}{北魏·昙曜主持开凿云冈早期大像，皇室与僧团共同塑造佛教王权象征}
  \item[460年] \eventanchor{NSL-440-589-037}{北魏·北魏在平城组织石工、供养人与物资参与云冈营造}
  \item[461年] \eventanchor{NSL-440-589-038}{北魏·文成朝对北部边镇维持军马与粮食调配}
  \item[462年] \eventanchor{NSL-440-589-039}{刘宋·宋廷在淮北方向调整守将与屯戍，避免再作大规模北伐}
  \item[463年] \eventanchor{NSL-440-589-040}{北魏·北魏继续遣使与河西、西域绿洲政权保持朝贡和交通联系}
  \item[464年] \eventanchor{NSL-440-589-041}{刘宋·孝武帝去世，太子刘子业即位}
  \item[465年] \eventanchor{NSL-440-589-042}{刘宋·前废帝刘子业以诛杀、羞辱宗室和大臣的方式集权，引发朝内恐惧}
  \item[465年] \eventanchor{NSL-440-589-043}{北魏·献文帝即位，冯太后继续掌握朝廷决策}
  \item[466年] \eventanchor{NSL-440-589-044}{刘宋·湘东王刘彧等发动政变，杀前废帝并立明帝}
  \item[466年] \eventanchor{NSL-440-589-045}{刘宋·晋安王刘子勋在寻阳称帝，江州、荆州等地响应}
  \item[466年] \eventanchor{NSL-440-589-046}{北魏与刘宋·北魏利用刘宋内战夺取淮北若干城镇并介入边境局势}
  \item[467年] \eventanchor{NSL-440-589-047}{刘宋·明帝军击败刘子勋集团，寻阳政权瓦解}
  \item[467年] \eventanchor{NSL-440-589-048}{刘宋·子勋之乱后南方州郡重新编置守将，流民与逃户问题凸显}
  \item[467年] \eventanchor{NSL-440-589-049}{南朝道教·陆修静整理灵宝、上清等道经目录与仪式传统的活动延续至此期}
  \item[468年] \eventanchor{NSL-440-589-050}{北魏·冯太后整顿朝廷人事，提升汉人士族与鲜卑贵族的协作}
  \item[469年] \eventanchor{NSL-440-589-051}{北魏·北魏延续云冈洞窟营造，题记和造像反映多层供养关系}
  \item[470年] \eventanchor{NSL-440-589-052}{刘宋·明帝继续限制宗室兵权并依赖近臣、禁军维持建康控制}
  \item[471年] \eventanchor{NSL-440-589-053}{北魏·献文帝禅位于太子拓跋宏，是为孝文帝}
  \item[471年] \eventanchor{NSL-440-589-054}{北魏·献文帝退居别宫后仍参与政事，形成非常规的皇权双重结构}
  \item[472年] \eventanchor{NSL-440-589-055}{北魏·献文帝去世，孝文帝与冯太后之间的摄政关系更为明确}
  \item[472年] \eventanchor{NSL-440-589-056}{刘宋·明帝加强对内廷和宗室的控制，朝政日益倚重近侍}
  \item[473年] \eventanchor{NSL-440-589-057}{北魏·北魏在平城朝廷中继续讨论俸禄、官吏考课与地方治理问题}
  \item[474年] \eventanchor{NSL-440-589-058}{刘宋·桂阳王刘休范自江州起兵，进逼建康后被平定}
  \item[474年] \eventanchor{NSL-440-589-059}{北魏·北魏继续以南部州镇防备宋军与接纳江淮来降者}
  \item[475年] \eventanchor{NSL-440-589-060}{刘宋·明帝去世，幼主刘昱即位，萧道成掌握禁军要职}
  \item[476年] \eventanchor{NSL-440-589-061}{刘宋·后废帝刘昱暴虐失序，禁军与大臣对其统治的容忍度下降}
  \item[477年] \eventanchor{NSL-440-589-062}{刘宋·萧道成杀后废帝，立顺帝并以相国身份控制朝政}
  \item[477年] \eventanchor{NSL-440-589-063}{刘宋·袁粲、沈攸之等反对萧道成的行动相继失败}
  \item[478年] \eventanchor{NSL-440-589-064}{刘宋·萧道成继续清除宋室支持者并重组州镇任命}
  \item[479年] \eventanchor{NSL-440-589-065}{南齐·萧道成受禅，建南齐，刘宋灭亡}
  \item[479年] \eventanchor{NSL-440-589-066}{北魏与南齐·北魏观察南朝改朝并调整淮北边防的对齐策略}
  \item[480年] \eventanchor{NSL-440-589-067}{南齐·齐高帝以赦令、封赏和州镇任命巩固新朝}
  \item[481年] \eventanchor{NSL-440-589-068}{北魏·北魏在冯太后主导下继续整饬官吏考课与俸禄议题}
  \item[482年] \eventanchor{NSL-440-589-069}{南齐·齐武帝即位，南齐进入相对稳定的永明初期}
  \item[483年] \eventanchor{NSL-440-589-070}{北魏·北魏筹议班禄与官吏供给，试图改变无俸官制下的掠夺性征收}
  \item[484年] \eventanchor{NSL-440-589-071}{北魏·北魏颁行班禄制，以谷帛等供给官员并配合考课}
  \item[484年] \eventanchor{NSL-440-589-072}{南齐·齐武帝经营建康朝廷并以儒学、礼制与选举巩固皇权}
  \item[485年] \eventanchor{NSL-440-589-074}{北魏·均田令把露田、桑田等类别与还受、继承规则相联系}
  \item[486年] \eventanchor{NSL-440-589-076}{北魏·三长制在部分地区与旧有宗族、坞壁和地方势力发生磨合}
  \item[487年] \eventanchor{NSL-440-589-077}{北魏·北魏继续按均田、三长框架整顿户籍和赋役}
  \item[488年] \eventanchor{NSL-440-589-078}{南齐·齐武帝在位时期继续维持淮南防线与北魏使节往来}
  \item[489年] \eventanchor{NSL-440-589-079}{北魏·冯太后去世，孝文帝开始亲政}
  \item[489年] \eventanchor{NSL-440-589-080}{北魏·冯太后陵寝与丧礼安排体现鲜卑皇室、汉式礼制与政治记忆的交织}
  \item[490年] \eventanchor{NSL-440-589-081}{北魏·孝文帝亲政后整顿中枢任官并压缩旧贵族的独立空间}
  \item[491年] \eventanchor{NSL-440-589-082}{南齐·齐武帝继续经营永明政治，士族文人和州郡官僚活跃于建康网络}
  \item[492年] \eventanchor{NSL-440-589-083}{南朝道教·陶弘景辞官入句曲山修道，连接士族、道教与医药知识网络}
  \item[493年] \eventanchor{NSL-440-589-084}{北魏·孝文帝以南伐为名率众南下，实际启动自平城迁都洛阳的战略转向}
  \item[493年] \eventanchor{NSL-440-589-085}{北魏·北魏朝臣围绕迁都利弊争论，皇帝以军事行动压制反对}
  \item[494年] \eventanchor{NSL-440-589-087}{北魏·洛阳城郭、宫城、道路和坊市的修建吸纳官营资源、工匠与迁入人口}
  \item[495年] \eventanchor{NSL-440-589-088}{北魏·孝文帝改易服饰、语言和礼仪，要求朝臣在洛阳遵行新制}
  \item[495年] \eventanchor{NSL-440-589-089}{北魏·北魏推动鲜卑贵族与汉人士族通婚、联姻及谱系整合}
  \item[496年] \eventanchor{NSL-440-589-090}{北魏·孝文帝改拓跋氏为元氏，并令部分贵族改用汉姓}
  \item[496年] \eventanchor{NSL-440-589-091}{北魏·北魏重定门第与官品秩序，试图以谱系、官爵组织新都精英}
  \item[497年] \eventanchor{NSL-440-589-092}{北魏与南齐·南齐将领王肃降魏，北魏借其熟悉南朝制度和地理参与南伐谋划}
  \item[497年] \eventanchor{NSL-440-589-093}{北魏·孝文帝以洛阳为基地筹划南伐，整顿南部州郡转运与军队}
  \item[498年] \eventanchor{NSL-440-589-094}{北魏与南齐·北魏大举南伐，围攻南阳、赭阳等淮汉要地}
  \item[498年] \eventanchor{NSL-440-589-095}{南齐·南齐在北魏南伐压力下调整荆、雍、豫州防务}
  \item[499年] \eventanchor{NSL-440-589-096}{北魏·孝文帝在南征途中去世，宣武帝即位}
  \item[499年] \eventanchor{NSL-440-589-097}{南齐·齐明帝去世，萧宝卷即位，宫廷政治迅速失序}
  \item[500年] \eventanchor{NSL-440-589-098}{北魏·宣武帝即位后依靠宗室与辅臣维持洛阳新都秩序}
  \item[500年] \eventanchor{NSL-440-589-099}{北魏·裴叔业以寿阳降魏，北魏取得淮南重要据点}
  \item[500年] \eventanchor{NSL-440-589-100}{南齐·陈显达等在南齐末政中起兵反萧宝卷，旋遭镇压}
  \item[501年] \eventanchor{NSL-440-589-101}{北魏·宣武帝镇压咸阳王禧叛乱，限制宗室在洛阳的独立军政资源}
  \item[501年] \eventanchor{NSL-440-589-102}{南齐·萧衍自襄阳起兵，联合雍州与荆州军政资源东下}
  \item[501年] \eventanchor{NSL-440-589-103}{南齐·萧宝卷被杀，萧衍拥立萧宝融于江陵并控制建康}
  \item[502年] \eventanchor{NSL-440-589-104}{梁·萧衍受禅建立梁，改元天监}
  \item[502年] \eventanchor{NSL-440-589-105}{梁·梁武帝发布宽刑、整饬吏治等建国初政以重建朝廷信誉}
  \item[502年] \eventanchor{NSL-440-589-106}{北魏与梁·北魏以寿阳等前沿观察梁朝建国后的边防与使节安排}
  \item[503年] \eventanchor{NSL-440-589-107}{梁·梁朝整编齐末降附军队与州郡官员，控制长江中下游漕运}
  \item[503年] \eventanchor{NSL-440-589-108}{北魏·北魏以寿阳为基地经营淮南屯戍和转运}
  \item[504年] \eventanchor{NSL-440-589-109}{梁·梁武帝以礼法、赦宥与官员任用区别于齐末失政}
  \item[504年] \eventanchor{NSL-440-589-110}{北魏·宣武朝继续任用洛阳士族和北方旧贵族，调和迁都后的官僚竞争}
  \item[505年] \eventanchor{NSL-440-589-111}{北魏与梁·北魏以元英等进攻梁的义阳、汉水防线}
  \item[505年] \eventanchor{NSL-440-589-112}{梁·梁将韦睿等组织义阳救援与淮汉防务}
  \item[506年] \eventanchor{NSL-440-589-113}{北魏与梁·钟离之战中梁军据淮水防御，北魏军撤退}
  \item[506年] \eventanchor{NSL-440-589-114}{梁·梁军利用钟离胜利整修淮南城戍与水运联络}
  \item[507年] \eventanchor{NSL-440-589-115}{北魏·宣武帝裁抑部分宗室与权臣，洛阳中枢围绕近侍和外戚展开竞争}
  \item[507年] \eventanchor{NSL-440-589-116}{梁·梁廷以建康为中心整理户籍、漕运与州郡赋役文书}
  \item[508年] \eventanchor{NSL-440-589-117}{北魏·北魏在洛阳周边继续扩展寺院、官署和贵族居住空间}
  \item[509年] \eventanchor{NSL-440-589-118}{北魏·北魏以南部州镇安置梁朝来降者并利用其熟悉江淮地理}
  \item[510年] \eventanchor{NSL-440-589-119}{北魏·宣武朝重视洛阳宫城、道路和水利供给，巩固新都日常运作}
  \item[510年] \eventanchor{NSL-440-589-120}{梁·梁武帝持续扶持寺院与讲经活动，佛教成为朝廷礼制之外的重要公共网络}
  \item[511年] \eventanchor{NSL-440-589-121}{北魏·北魏南部守军与梁军在淮汉一线维持小规模攻守和侦察}
  \item[512年] \eventanchor{NSL-440-589-122}{梁·梁朝调整荆、郢、豫诸州守将，防范北魏与地方军府的双重压力}
  \item[513年] \eventanchor{NSL-440-589-123}{北魏·宣武帝晚年近臣与后族参与中枢决策，外廷制衡能力下降}
  \item[514年] \eventanchor{NSL-440-589-124}{南朝佛教·僧祐编纂《弘明集》，汇集佛教与儒道辩难文献}
  \item[515年] \eventanchor{NSL-440-589-125}{北魏·宣武帝去世，幼主孝明帝即位，胡太后临朝}
  \item[515年] \eventanchor{NSL-440-589-126}{北魏·胡太后复兴佛教营建并以施舍、造像建立摄政权威}
  \item[515年] \eventanchor{NSL-440-589-127}{南朝佛教·僧祐完成《出三藏记集》编纂活动，梳理译经、经录与僧传材料}
  \item[516年] \eventanchor{NSL-440-589-128}{北魏·胡太后下令营建永宁寺及高塔，形成洛阳突出的宗教地标}
  \item[516年] \eventanchor{NSL-440-589-129}{北魏·永宁寺营建动员木材、工匠和周边运输，反映洛阳城市供给体系}
  \item[517年] \eventanchor{NSL-440-589-130}{北魏·元叉发动政变，幽禁胡太后并控制孝明帝与朝政}
  \item[517年] \eventanchor{NSL-440-589-131}{梁·梁朝持续与北魏互遣使节，围绕边境、礼仪与降人问题交涉}
  \item[518年] \eventanchor{NSL-440-589-132}{北魏·宋云、惠生受命西行求法，途经西域并记录佛教与交通见闻}
  \item[518年] \eventanchor{NSL-440-589-133}{北魏·元叉集团继续控制官员任免，地方州镇对洛阳政治的信任减弱}
  \item[519年] \eventanchor{NSL-440-589-134}{北魏·胡太后复出并诛元叉，重新临朝}
  \item[520年] \eventanchor{NSL-440-589-135}{北魏·孝明朝整修洛阳寺院、官署与都城供给体系}
  \item[约520年] \eventanchor{NSL-440-589-136}{北魏与南朝佛教·菩提达摩入华和与梁武帝会见的叙事在后世禅宗传记中形成}
  \item[521年] \eventanchor{NSL-440-589-137}{梁·梁武帝继续组织讲经、戒律与僧官相关活动，加强朝廷与僧团联系}
  \item[522年] \eventanchor{NSL-440-589-138}{北魏·北魏北部六镇军民的待遇、迁徙和仕进问题日益尖锐}
  \item[523年] \eventanchor{NSL-440-589-139}{北魏·沃野镇破六韩拔陵起兵，六镇之乱爆发}
  \item[523年] \eventanchor{NSL-440-589-140}{北魏·怀朔、武川等镇相继响应或受叛乱波及}
  \item[523年] \eventanchor{NSL-440-589-141}{北魏·六镇动乱扰乱北部牧地、军粮和户籍，流民向并、肆、河北移动}
  \item[524年] \eventanchor{NSL-440-589-142}{北魏·柔玄镇杜洛周起兵，河北北部叛乱扩大}
  \item[524年] \eventanchor{NSL-440-589-143}{北魏·破六韩拔陵军与官军反复争战，北魏求助柔然力量}
  \item[524年] \eventanchor{NSL-440-589-144}{北魏·孝明朝在叛乱中增派州郡征发，河北、并州地方负担上升}
  \item[525年] \eventanchor{NSL-440-589-145}{北魏·北魏在六镇战争中分化叛军并安置部分降附者}
  \item[525年] \eventanchor{NSL-440-589-146}{北魏·北方叛乱使洛阳与地方间的漕运、市场和道路安全恶化}
  \item[526年] \eventanchor{NSL-440-589-147}{北魏·尔朱荣以并州军镇力量受命镇压北方叛乱}
  \item[526年] \eventanchor{NSL-440-589-148}{北魏·尔朱荣部在北方作战并吸纳降卒、部曲和流民}
  \item[527年] \eventanchor{NSL-440-589-149}{梁·梁武帝首次舍身同泰寺，由朝廷以赎金迎还}
  \item[527年] \eventanchor{NSL-440-589-150}{梁·同泰寺及建康寺院活动带动斋会、施舍与城市宗教聚集}
  \item[527年] \eventanchor{NSL-440-589-151}{北魏·萧宝夤、万俟丑奴等关陇叛乱集团与北魏军反复交战}
  \item[528年] \eventanchor{NSL-440-589-152}{北魏·孝明帝被杀，胡太后先立女婴后改立元钊，引发继承危机}
  \item[528年] \eventanchor{NSL-440-589-153}{北魏·尔朱荣入洛阳，河阴大规模屠杀朝臣与士人}
  \item[528年] \eventanchor{NSL-440-589-154}{北魏·尔朱荣立孝庄帝，北魏皇权受制于并州军人集团}
  \item[528年] \eventanchor{NSL-440-589-155}{北魏·北方战乱与河阴屠杀促使士人、僧侣和民众外迁或避乱}
  \item[529年] \eventanchor{NSL-440-589-156}{北魏·梁朝扶植的元颢渡淮入洛阳，短暂夺取北魏都城}
  \item[529年] \eventanchor{NSL-440-589-157}{北魏·尔朱荣自北方回师击败元颢，恢复孝庄帝于洛阳}
  \item[529年] \eventanchor{NSL-440-589-158}{梁·梁朝对元颢行动的投入未能转化为持续北方控制}
  \item[530年] \eventanchor{NSL-440-589-159}{北魏·孝庄帝设计杀死尔朱荣，试图恢复皇帝自主权}
  \item[530年] \eventanchor{NSL-440-589-160}{北魏·尔朱兆等攻入洛阳，杀孝庄帝并重建尔朱集团统治}
  \item[531年] \eventanchor{NSL-440-589-161}{北魏·高欢在信都起兵反尔朱氏，联合河北豪强和流民武装}
  \item[531年] \eventanchor{NSL-440-589-162}{北魏·高欢集团利用河北粮仓、坞壁和地方精英建立供给网络}
  \item[532年] \eventanchor{NSL-440-589-163}{北魏·高欢在韩陵击败尔朱氏主力，控制洛阳政治}
  \item[532年] \eventanchor{NSL-440-589-164}{北魏·高欢立孝武帝，安排亲信掌握军政要职}
  \item[533年] \eventanchor{NSL-440-589-165}{北魏·高欢以晋阳、邺城和河北州郡为基地组织军政}
  \item[533年] \eventanchor{NSL-440-589-166}{北魏·孝武帝试图联络贺拔胜、宇文泰等西部军人以制衡高欢}
  \item[534年] \eventanchor{NSL-440-589-168}{西魏·宇文泰迎孝武帝入长安，掌握关中军政}
  \item[534年] \eventanchor{NSL-440-589-169}{东魏·高欢迁孝静帝与百官至邺，建立东部政权中心}
  \item[535年] \eventanchor{NSL-440-589-170}{西魏·孝武帝死后，宇文泰立元宝炬为文帝，西魏正式建国}
  \item[535年] \eventanchor{NSL-440-589-171}{东魏·孝静帝在邺即位，东魏正式建国}
  \item[536年] \eventanchor{NSL-440-589-172}{西魏·宇文泰整顿关中军队与州郡，吸纳流民和地方豪强}
  \item[536年] \eventanchor{NSL-440-589-173}{东魏·高欢在邺调整河北官员与军队供给，巩固东魏新政权}
  \item[537年] \eventanchor{NSL-440-589-174}{东魏与西魏·宇文泰在沙苑击败高欢军}
  \item[537年] \eventanchor{NSL-440-589-175}{西魏·沙苑战后西魏安置俘虏与战时流民，修补关中供给}
  \item[538年] \eventanchor{NSL-440-589-176}{东魏与西魏·河桥、洛阳一线再成双方争夺战场，战局反复}
  \item[538年] \eventanchor{NSL-440-589-177}{西魏·宇文泰经营关中屯田、仓储与军户供给以承受东魏压力}
  \item[539年] \eventanchor{NSL-440-589-178}{东魏与西魏·东西魏一度议和并交换使节，未能消除河南、关中争端}
  \item[540年] \eventanchor{NSL-440-589-179}{东魏·高欢继续以晋阳为北方军事重镇，调动河北资源支援前线}
  \item[541年] \eventanchor{NSL-440-589-180}{梁与交州·交州豪族李贲起兵反梁，地方治理危机显现}
  \item[542年] \eventanchor{NSL-440-589-181}{梁与交州·梁军讨伐李贲，交州战争与气候、道路和补给困难相互影响}
  \item[543年] \eventanchor{NSL-440-589-182}{东魏与西魏·邙山之战中高欢与宇文泰激战，双方均未取得决定性歼灭}
  \item[544年] \eventanchor{NSL-440-589-183}{万春国·李贲在交州称帝，建立万春国的政治名号}
  \item[544年] \eventanchor{NSL-440-589-184}{西魏·宇文泰持续以关中军户和部曲补充军队，整合地方豪强}
  \item[545年] \eventanchor{NSL-440-589-185}{梁·徐陵主持或参与《玉台新咏》的编纂活动，保存宫体诗与女性题材文本}
  \item[545年] \eventanchor{NSL-440-589-186}{东魏·高欢继续通过邺、晋阳两地调度军政，稳定河北内部}
  \item[546年] \eventanchor{NSL-440-589-187}{东魏·高欢去世，高澄继承东魏军政主导权}
  \item[547年] \eventanchor{NSL-440-589-188}{东魏与梁·侯景叛高澄后降梁，带领部众进入淮南}
  \item[547年] \eventanchor{NSL-440-589-189}{东魏·杨衒之在东魏时期撰成《洛阳伽蓝记》，追述北魏洛阳寺院与毁乱}
  \item[547年] \eventanchor{NSL-440-589-190}{梁·梁朝围绕安置侯景军队、粮饷和驻防地点与其发生冲突}
  \item[548年] \eventanchor{NSL-440-589-191}{梁·侯景举兵反梁，渡江进攻建康}
  \item[548年] \eventanchor{NSL-440-589-192}{梁·侯景围建康，城内粮食、疫病与军民秩序恶化}
  \item[549年] \eventanchor{NSL-440-589-193}{梁·建康陷落，梁武帝被侯景控制并去世}
  \item[549年] \eventanchor{NSL-440-589-194}{梁·侯景以萧正德、后又以萧纲等为傀儡操控朝廷}
  \item[550年] \eventanchor{NSL-440-589-196}{北齐·高洋废东魏孝静帝，自立为帝，北齐建立}
  \item[550年] \eventanchor{NSL-440-589-197}{北齐·北齐接收东魏旧官僚、军户与州郡文书，重整河北统治}
  \item[550年] \eventanchor{NSL-440-589-198}{北齐·高洋即位后加强对鲜卑、汉人军政集团的赏罚与编制}
  \item[550年] \eventanchor{NSL-440-589-199}{梁·侯景废简文帝并自立为汉帝，继续控制建康}
  \item[551年] \eventanchor{NSL-440-589-200}{梁·侯景军与湘东王萧绎、王僧辩等势力在江汉、江东对抗}
  \item[551年] \eventanchor{NSL-440-589-201}{北齐·北齐与柔然、突厥等草原势力就边境和婚姻关系展开交涉}
  \item[552年] \eventanchor{NSL-440-589-202}{突厥与柔然·突厥首领土门击败柔然可汗，突厥汗国兴起}
  \item[552年] \eventanchor{NSL-440-589-203}{突厥·土门死后，突厥汗位由科罗等继承，汗国继续向东西扩展}
  \item[552年] \eventanchor{NSL-440-589-204}{梁·王僧辩、陈霸先等在江南击败侯景主力，收复建康}
  \item[552年] \eventanchor{NSL-440-589-205}{梁·侯景之乱后建康、京口和江州的户籍、仓储与寺院需重新恢复}
  \item[553年] \eventanchor{NSL-440-589-206}{梁·侯景逃亡途中被部下杀害，叛乱核心集团瓦解}
  \item[553年] \eventanchor{NSL-440-589-207}{北齐·北齐与南朝围绕淮南、江北流民和边将展开安置与交涉}
  \item[553年] \eventanchor{NSL-440-589-208}{突厥与柔然·突厥继续攻击柔然残部，部分柔然人向西魏、北齐边地求援或流散}
  \item[554年] \eventanchor{NSL-440-589-209}{梁与西魏·西魏攻陷江陵，杀梁元帝萧绎并掳掠人口物资}
  \item[554年] \eventanchor{NSL-440-589-210}{西魏·西魏在江陵战后扶植萧詧，建立后梁的政治基础}
  \item[554年] \eventanchor{NSL-440-589-211}{北齐·北齐利用梁朝江陵失陷的局势介入江淮与南朝宗室争夺}
  \item[554年] \eventanchor{NSL-440-589-212}{北齐·高洋去世，高殷即位，杨愔等辅政}
  \item[555年] \eventanchor{NSL-440-589-213}{北齐·高演发动政变废高殷，自立为孝昭帝}
  \item[555年] \eventanchor{NSL-440-589-214}{梁·王僧辩迎北齐扶持的萧渊明入建康，试图重建梁朝}
  \item[555年] \eventanchor{NSL-440-589-215}{梁·陈霸先杀王僧辩，控制建康并改立萧方智}
  \item[555年] \eventanchor{NSL-440-589-216}{北齐与梁·北齐军南下支持萧渊明，后因江南抵抗与补给困难撤退}
  \item[556年] \eventanchor{NSL-440-589-217}{西魏·宇文泰去世，宇文护接掌西魏军政}
  \item[556年] \eventanchor{NSL-440-589-218}{北齐·高演去世，高湛即位为武成帝}
  \item[556年] \eventanchor{NSL-440-589-219}{梁·陈霸先整合京口、建康和江东军队，压制梁末宗室竞争者}
  \item[557年] \eventanchor{NSL-440-589-220}{北周·宇文护废西魏恭帝，立宇文觉，北周建立}
  \item[557年] \eventanchor{NSL-440-589-221}{北周·北周仿《周礼》设置六官，重构中央官制名目}
  \item[557年] \eventanchor{NSL-440-589-222}{陈·陈霸先受禅建立陈，梁朝终结}
  \item[557年] \eventanchor{NSL-440-589-223}{陈·陈朝接收梁末州郡文书、军人和侨民，重整建康周边秩序}
  \item[558年] \eventanchor{NSL-440-589-224}{北周·宇文护废宇文觉并立宇文毓为明帝，显示辅政者控制皇统}
  \item[558年] \eventanchor{NSL-440-589-225}{陈·陈文帝即位后继续压制王琳等长江中游势力}
  \item[559年] \eventanchor{NSL-440-589-226}{陈·王琳在长江中游拥立梁永嘉王萧庄并联络北齐}
  \item[559年] \eventanchor{NSL-440-589-227}{北齐·北齐以支持萧庄、王琳的方式介入长江中游局势}
  \item[560年] \eventanchor{NSL-440-589-228}{北齐·高湛禅位于高纬，北齐后主即位，武成帝仍掌实权}
  \item[560年] \eventanchor{NSL-440-589-229}{陈·陈文帝继续以江州、湘州等军镇防御王琳与北齐压力}
  \item[561年] \eventanchor{NSL-440-589-230}{北周与突厥·北周与突厥围绕和亲、贡赐和共同对齐的战略开展谈判}
  \item[561年] \eventanchor{NSL-440-589-231}{北周·宇文护继续掌握军队与中枢任命，北周皇帝难以自主}
  \item[562年] \eventanchor{NSL-440-589-232}{陈·陈军击败王琳集团，萧庄与王琳势力北走}
  \item[562年] \eventanchor{NSL-440-589-233}{陈·王琳部众与随军人口北迁，改变江汉、淮南的流民与军户分布}
  \item[563年] \eventanchor{NSL-440-589-234}{北齐·北齐在河清年间整饬律令、官制与州郡行政}
  \item[563年] \eventanchor{NSL-440-589-235}{北周·北周与北齐在河东、洛阳方向继续争夺城戍和交通}
  \item[564年] \eventanchor{NSL-440-589-236}{北周与北齐·北周大举攻北齐，围攻洛阳一线而未获决定性成果}
  \item[564年] \eventanchor{NSL-440-589-237}{北齐·北齐在北周进攻中利用河南城防、河北后勤组织反击}
  \item[565年] \eventanchor{NSL-440-589-238}{突厥与北齐·突厥对北齐边境施压并在北周、北齐之间调整盟约}
  \item[565年] \eventanchor{NSL-440-589-239}{北周·北周以关中军户、府兵和州郡粮运维持对齐备战}
  \item[566年] \eventanchor{NSL-440-589-240}{陈·陈文帝病逝，废帝即位，朝廷由重臣辅政}
  \item[567年] \eventanchor{NSL-440-589-241}{陈·陈顼控制朝政并处理建康宗室、军将与地方任命}
  \item[568年] \eventanchor{NSL-440-589-242}{陈·陈顼废陈伯宗为临海王，独揽朝政}
  \item[569年] \eventanchor{NSL-440-589-243}{陈·陈顼即皇帝位，是为宣帝}
  \item[569年] \eventanchor{NSL-440-589-244}{北齐·北齐后主朝内廷近臣、后族与宗室竞争加剧}
  \item[570年] \eventanchor{NSL-440-589-245}{北周·北周继续调整府兵统领与关中州郡资源，准备再攻北齐}
  \item[571年] \eventanchor{NSL-440-589-246}{北齐·北齐后主诛杀权臣和士开，宫廷派系重新洗牌}
  \item[572年] \eventanchor{NSL-440-589-247}{北周·周武帝诛杀宇文护，开始亲政}
  \item[572年] \eventanchor{NSL-440-589-248}{北周·周武帝整顿军队、赏罚与官员任用，削弱宇文护旧党}
  \item[573年] \eventanchor{NSL-440-589-249}{陈·陈宣帝北伐，吴明彻等攻取淮南、寿阳附近多处城镇}
  \item[573年] \eventanchor{NSL-440-589-250}{陈·陈朝在新取淮南地区安排守将、户籍和粮运}
  \item[573年] \eventanchor{NSL-440-589-251}{北齐·北齐在陈军北伐压力下调整淮南守备与河北援军}
  \item[574年] \eventanchor{NSL-440-589-252}{北周·周武帝下令禁断佛教、道教，毁寺并令僧道还俗}
  \item[574年] \eventanchor{NSL-440-589-253}{北周·北周禁佛伴随寺产、僧尼与宗教器物的清查和再分配}
  \item[574年] \eventanchor{NSL-440-589-254}{北齐·北齐以佛教护持区别于北周的禁佛政策，寺院和造像继续获得赞助}
  \item[575年] \eventanchor{NSL-440-589-255}{北周与北齐·周武帝亲率伐齐，进攻河东、河南方向}
  \item[575年] \eventanchor{NSL-440-589-256}{北齐·北齐组织河北、并州军队抵御北周，依赖城防与快速调兵}
  \item[576年] \eventanchor{NSL-440-589-257}{北周与北齐·北周攻克平阳，打开进取晋阳与邺城的通道}
  \item[576年] \eventanchor{NSL-440-589-258}{北齐·北齐后主在晋阳、邺城间调兵失当，军心与将领信任下降}
  \item[576年] \eventanchor{NSL-440-589-259}{北周·北周军在占领区接收城镇、粮仓与降附军民}
  \item[577年] \eventanchor{NSL-440-589-261}{北周·北周处置北齐皇室、官员、军户和寺院财产}
  \item[577年] \eventanchor{NSL-440-589-262}{北周·北周在原北齐地区恢复佛教活动的幅度与禁令执行出现差异}
  \item[578年] \eventanchor{NSL-440-589-263}{北周·周武帝去世，周宣帝即位}
  \item[578年] \eventanchor{NSL-440-589-264}{北周与突厥·突厥趁北周皇位更替施加边境压力，北周需维持和亲与防务}
  \item[579年] \eventanchor{NSL-440-589-265}{北周·周宣帝禅位于幼子静帝，自为天元皇帝}
  \item[579年] \eventanchor{NSL-440-589-266}{北周·北周在河北推行州郡整并、户籍和赋役恢复措施}
  \item[579年] \eventanchor{NSL-440-589-267}{北周·北周对佛教禁令有所松动，部分寺院和僧侣重新活动}
  \item[580年] \eventanchor{NSL-440-589-268}{北周·周宣帝去世，杨坚以外戚身份辅政静帝}
  \item[580年] \eventanchor{NSL-440-589-269}{北周·尉迟迥在相州起兵反杨坚，河北旧齐势力与部分军将响应}
  \item[580年] \eventanchor{NSL-440-589-270}{北周·杨坚平定尉迟迥后重新任命河北州郡和军队将领}
  \item[581年] \eventanchor{NSL-440-589-272}{隋·隋初重整中央官制、府兵与州郡任命，削减北周宗室影响}
  \item[581年] \eventanchor{NSL-440-589-273}{隋与突厥·突厥沙钵略可汗支持北周宗室反隋，北部边境紧张}
  \item[582年] \eventanchor{NSL-440-589-274}{陈·陈宣帝去世，陈叔宝即位，是为后主}
  \item[582年] \eventanchor{NSL-440-589-275}{隋·隋营建大兴城，规划宫城、官署、坊市与交通供给}
  \item[582年] \eventanchor{NSL-440-589-276}{隋与突厥·突厥大举南下，隋以长城、州镇和机动军队应对}
  \item[583年] \eventanchor{NSL-440-589-277}{隋·隋修订开皇律令，强调简约刑法与中央审断}
  \item[583年] \eventanchor{NSL-440-589-278}{隋·隋整并州郡、清理户籍与赋役，推进均田和租调体系}
  \item[583年] \eventanchor{NSL-440-589-279}{隋与突厥·隋利用突厥可汗间竞争，削弱沙钵略对北周遗臣的支持}
  \item[584年] \eventanchor{NSL-440-589-280}{隋·隋开通广通渠，连接大兴城与潼关、黄河运输体系}
  \item[584年] \eventanchor{NSL-440-589-281}{陈·陈朝面对隋北方整合，修整长江防线并调整江北守将}
  \item[585年] \eventanchor{NSL-440-589-282}{隋·隋在北方旧齐地区继续编户、核田与恢复仓储转运}
  \item[585年] \eventanchor{NSL-440-589-283}{陈·陈后主朝廷的奢侈、宫廷文学与近侍政治受到史家强烈批评}
  \item[586年] \eventanchor{NSL-440-589-284}{隋·隋继续整顿府兵、仓储和江淮方向的军情侦察}
  \item[586年] \eventanchor{NSL-440-589-285}{陈·陈朝在长江下游维持水军、城戍与漕运，承受持续军费}
  \item[587年] \eventanchor{NSL-440-589-286}{隋·隋废后梁，接收荆襄州郡与江陵资源}
  \item[587年] \eventanchor{NSL-440-589-287}{隋·隋在荆襄安置后梁官员、军人和居民，整顿道路与粮运}
  \item[588年] \eventanchor{NSL-440-589-288}{隋·隋以杨广、杨俊、杨素等分路部署，发动灭陈战争}
  \item[588年] \eventanchor{NSL-440-589-289}{陈·陈朝调集长江水军和沿岸守将抵御隋军}
  \item[589年] \eventanchor{NSL-440-589-291}{隋·隋接收陈朝州郡、户籍、仓储和官员，派使安抚江南}
  \item[589年] \eventanchor{NSL-440-589-292}{隋·隋处理陈室宗族、降将与江南士族的封授和迁置}
  \item[约589年] \eventanchor{NSL-440-589-293}{北齐遗民与隋·颜之推晚年完成或修订《颜氏家训》的主要篇章，反思北齐、北周、隋之间的家族处境}
  \item[约550年] \eventanchor{NSL-440-589-294}{北齐·北齐在邺城重整东魏旧都的宫署、市场与军政居住空间}
  \item[552年] \eventanchor{NSL-440-589-295}{北齐与柔然残部·北齐在北部边地处置柔然覆亡后逃入境内的部众}
  \item[554年] \eventanchor{NSL-440-589-296}{西魏与梁·江陵陷落后，大批梁朝官员、工匠和图书被迁徙或散失}
  \item[557年] \eventanchor{NSL-440-589-297}{北周·北周建国初年以六官名目协调军政、财政与礼制机构}
  \item[559年] \eventanchor{NSL-440-589-298}{陈·陈朝在建康战后修复宫城、仓储和长江水运节点}
  \item[563年] \eventanchor{NSL-440-589-299}{北齐·北齐河清制度整理涉及律令、官品与地方行政的重新表述}
  \item[566年] \eventanchor{NSL-440-589-300}{北齐与突厥·北齐继续以贡赐和使节维系对突厥的边境关系}
  \item[570年] \eventanchor{NSL-440-589-301}{北周与北齐·北周侦察并经营河东、河南道路，准备对北齐发动更大攻势}
  \item[575年] \eventanchor{NSL-440-589-302}{北周·禁佛后的北周以州郡接收寺院人户与土地，重估其赋役归属}
  \item[577年] \eventanchor{NSL-440-589-303}{北周·北周在新并北齐地区重建仓储、道路与军队粮运}
  \item[580年] \eventanchor{NSL-440-589-304}{隋前夜·尉迟迥之乱中，杨坚军利用关中与河北粮道协调平叛}
  \item[582年] \eventanchor{NSL-440-589-305}{隋·大兴城建设同步规划水源、道路与坊市，吸纳关中工匠和劳力}
\end{description}
